\section{Proofs for Lie-Algebra Links}
\label{sec:proofs-links}

\noindent\emph{Proof of \Cref{prop:link-chart-image}.}
(a) A link family lies in the image exactly when
$\mathring\Theta^x_e=U_i^{-1}\Theta^x_eU_j\in\exp(\mathfrak g)$ for every $e\in E_+$, and then for every
$e\in E$ because $\exp(\mathring\omega)^{-1}=\exp(-\mathring\omega)$. Passive re-expressions leave
twists unchanged, reference changes and active automorphisms conjugate them
(\Cref{prop:geo-edge-twist-invariance}), and $\exp(\mathfrak g)$ is closed under conjugation by
\Cref{lem:link-exp-naturality}.
(b) The first statement is \Cref{prop:link-real-logarithm}. Jacobi's formula
$\frac{\mathrm d}{\mathrm dt}\det e^{t\mathbf H}=\Tr(\mathbf H)\det e^{t\mathbf H}$ gives
$\det e^{\mathbf H}=e^{\Tr\mathbf H}>0$. If $\operatorname{diag}(-1,-2,1,\ldots,1)$ were $e^{\mathbf H}$
with $\mathbf H$ real, then $\mathbf H$ would commute with it, since it is a power series in
$\mathbf H$, and would therefore preserve its eigenspace for the eigenvalue $-1$, which is spanned by
the first standard basis vector $\mathbf b$; with $\mathbf H\mathbf b=\mathrm c\mathbf b$ for a real $\mathrm c$, this
gives $e^{\mathbf H}\mathbf b=e^{\mathrm c}\mathbf b\neq-\mathbf b$, a contradiction.
(c) For $\mathbf B\in\GL^+(K)$ the spectral theorem writes the positive definite matrix
$\mathbf B\mathbf B^\top$ as $\mathbf O_0\operatorname{diag}(d_1,\ldots,d_K)\mathbf O_0^\top$ with
$\mathbf O_0$ orthogonal and every $d_k>0$. The symmetric matrix
$\mathbf S=\mathbf O_0\operatorname{diag}(\frac12\log d_1,\ldots,\frac12\log d_K)\mathbf O_0^\top$ has
$e^{2\mathbf S}=\mathbf B\mathbf B^\top$, and $\mathbf O_1=e^{-\mathbf S}\mathbf B$ satisfies
$\mathbf O_1\mathbf O_1^\top=e^{-\mathbf S}\mathbf B\mathbf B^\top e^{-\mathbf S}=I_K$ and
$\det\mathbf O_1=\det\mathbf B/\det e^{\mathbf S}>0$, so $\mathbf O_1\in SO(K)$. By the real normal
form of an orthogonal matrix there is an orthogonal $\mathbf O_2$ such that
$\mathbf O_2^\top\mathbf O_1\mathbf O_2$ is block diagonal with diagonal entries $\pm1$ and $2\times2$
blocks $R(\phi)=\exp(\phi(E_{21}-E_{12}))$; since $\det\mathbf O_1=1$ the entries $-1$ occur an even
number of times and pair into blocks $R(\pi)$, so $\mathbf O_2^\top\mathbf O_1\mathbf O_2=\exp(\mathbf T_0)$
with $\mathbf T_0$ block diagonal and skew-symmetric. Hence $\mathbf O_1=\exp(\mathbf T)$ with
$\mathbf T=\mathbf O_2\mathbf T_0\mathbf O_2^\top$ skew-symmetric, $\mathbf B=\exp(\mathbf S)\exp(\mathbf T)$,
and the path $t\mapsto\exp(t\mathbf S)\exp(t\mathbf T)$ joins $I_K$ to $\mathbf B$ inside $\GL^+(K)$.
The Lie algebra of $SO(K)$ consists of the skew-symmetric matrices, so for $G=SO(K)$ every twist is an
exponential.
(d) The differential of $\exp$ at $0$ is the identity of $\mathfrak g$, since $\exp(t\mathring\omega)$
has velocity $\mathring\omega$ at $t=0$, so by the inverse function theorem $\exp(\mathfrak g)$
contains an open neighborhood of $e_G$. The subgroup generated by $\exp(\mathfrak g)$ is therefore a
union of translates of that neighborhood, hence open; its complement is a union of its cosets, hence
open; and a nonempty subset of a connected group that is open and closed is the whole group.
(e) Let the two pairs be traversed at the positions $a<a'$ of $p$. Given $\mathbf B\in\GL^+(K)$, write
$\mathbf B=\exp(\mathbf S)\exp(\mathbf T)$ by (c), give the pair traversed at $a$ the generator
$\pm\mathbf S$ and the pair traversed at $a'$ the generator $\pm\mathbf T$, with the signs chosen by
\cref{eq:link-generator-reversal} so that the traversed copies carry $\mathbf S$ and $\mathbf T$, and give
every other pair the generator $0$. Every other twist along $p$ is then $I_K$, and
$\mathring\Theta_p=\mathbf B$. The displayed case is
$\operatorname{diag}(1,2)R(\pi)=\operatorname{diag}(-1,-2)$, which is not an exponential by (b).
\hfill$\square$

\noindent\emph{Proof of \Cref{prop:link-loop-defect}.}
Write $\mathsf G=\mathsf G^{\mathrm{fib}}$ and $\Delta=I_{\mathbb V}-\rho_x(\mathring\Theta^x_p)$, and
note that, for invertible $\mathbf R$ and symmetric positive definite $\mathsf G'$, cyclicity of the trace
gives
\begin{equation}
\lVert\mathbf R^{-1}\Delta\mathbf R\rVert^2_{\mathsf G'}
=\Tr\bigl(\mathbf R\mathsf G'^{-1}\mathbf R^\top\Delta^\top\mathbf R^{-\top}\mathsf G'\mathbf R^{-1}\Delta\bigr)
=\lVert\Delta\rVert^2_{\mathbf R^{-\top}\mathsf G'\mathbf R^{-1}},
\label{eq:link-defect-conjugation}
\end{equation}
where a subscript names the matrix used in \cref{eq:link-trace-form}.

(a) Cyclicity gives
$\langle\mathbf H,\mathbf H'\rangle_{\mathsf G}=\Tr\bigl((\mathsf G^{1/2}\mathbf H\mathsf G^{-1/2})^\top(\mathsf G^{1/2}\mathbf H'\mathsf G^{-1/2})\bigr)$,
the Frobenius inner product of the images under the algebra isomorphism
$\mathbf H\mapsto\mathsf G^{1/2}\mathbf H\mathsf G^{-1/2}$, which yields the inner product,
submultiplicativity, and $\lVert I_{\mathbb V}\rVert_{\mathsf G}=\lVert I_{\mathbb V}\rVert_{\mathrm F}$. The
defect is continuous because $\exp$, products, and $\rho_x$ are; it depends only on the twists along
$p$, which are functions of the links; its zero set is that of a norm; and
$\rho_x(\mathrm{Hol}^x(p))=\rho_x(U_{i_0})\rho_x(\mathring\Theta^x_p)\rho_x(U_{i_0})^{-1}$ by
\cref{eq:link-holonomy-chart}.

(b) A passive re-expression changes neither the generators relative to $\sigma$
(\Cref{prop:link-generator-gauge}) nor $\sigma$ and $\mathsf G$. Relative to $\sigma'$ the reference
holonomy is $g_0^{-1}\mathring\Theta^x_pg_0$ (\Cref{prop:link-generator-gauge}), whose defect matrix is
$\mathbf R^{-1}\Delta\mathbf R$. The quotient convention \cref{eq:geo-quotient-convention} gives
$[\sigma,v]=[\sigma'\cdot g_0^{-1},v]=[\sigma',\mathbf R^{-1}v]$, so the coordinates of $\sigma'$ are
$\mathbf R^{-1}$ times those of $\sigma$, and the inner product with matrix $\mathsf G$ at $\sigma$ has
the matrix $\mathbf R^\top\mathsf G\mathbf R$ at $\sigma'$. Taking $\mathsf G'=\mathbf R^\top\mathsf G\mathbf R$
in \cref{eq:link-defect-conjugation} returns $\lVert\Delta\rVert^2_{\mathsf G}$, and $\mathsf G'=\mathsf G$
gives the second statement. If $\mathbf R^\top\mathsf G\mathbf R=\mathrm c_{\mathrm o}^2\mathsf G$ with
$\mathrm c_{\mathrm o}\neq0$, then
$\mathbf R^{-\top}\mathsf G\mathbf R^{-1}=\mathrm c_{\mathrm o}^{-2}\mathsf G$, and
\Cref{lem:link-trace-form-scale} gives equality. Conversely, let $p$ traverse the pair of $e$
exactly once, let $\mathbf H\in\mathfrak{gl}(\mathbb V)$, choose $\mathring\omega\in\mathfrak g$ with
$\mathrm d\rho_x\mathring\omega=\mathbf H$, and give the pair of $e$ the generator $t\mathring\omega$ and
every other pair the generator $0$. Then $\rho_x(\mathring\Theta^x_p)=\exp(\pm t\mathbf H)$ by
\Cref{lem:link-exp-naturality}, so $\Delta=\mp t\mathbf H+O(t^2)$. Dividing the assumed equality
$\lVert\Delta\rVert^2_{\mathbf R^{-\top}\mathsf G\mathbf R^{-1}}=\lVert\Delta\rVert^2_{\mathsf G}$ by $t^2$ and
letting $t\to0$ gives
$\lVert\mathbf H\rVert_{\mathbf R^{-\top}\mathsf G\mathbf R^{-1}}=\lVert\mathbf H\rVert_{\mathsf G}$ for every
$\mathbf H$, so \Cref{lem:link-trace-form-scale} gives
$\mathbf R^{-\top}\mathsf G\mathbf R^{-1}=\mathrm c_{\mathrm s}\mathsf G$ with $\mathrm c_{\mathrm s}>0$, and
$\mathrm c_{\mathrm s}^{1/2}\mathbf R$ is a $\mathsf G$-isometry. For the rotated walk,
$\mathring\Theta^x_{(e_1,\ldots,e_{r-1},e_0)}=(\mathring\Theta^x_{e_0})^{-1}\mathring\Theta^x_p\mathring\Theta^x_{e_0}$,
and \cref{eq:link-defect-conjugation} applies with $\mathbf R=\rho_x(\mathring\Theta^x_{e_0})$ and
$\mathsf G'=\mathsf G$.

(c) By the quotient convention, $[\sigma,v]=[u_i\cdot U_i,v]=[u_i,\rho_x(U_i)v]$, so the coordinates of
$u_i$ are $\rho_x(U_i)$ times those of $\sigma$, which gives the matrix $\mathsf G^{\mathrm{fib}}_i$. The
law under passive re-expression follows from $U_i'=g_iU_i$ in \cref{eq:geo-regime-one-gauge-law}, and
the relation between agents from $T_{ij}=U_iU_j^{-1}$. Since
$I_{\mathbb V}-\mathsf U=\rho_x(U_{i_0})\Delta\rho_x(U_{i_0})^{-1}$, \cref{eq:link-defect-conjugation} with
$\mathbf R=\rho_x(U_{i_0})^{-1}$ and $\mathsf G'=\mathsf G^{\mathrm{fib}}_{i_0}$ gives
$\lVert I_{\mathbb V}-\mathsf U\rVert^2_{\mathsf G^{\mathrm{fib}}_{i_0}}=\lVert\Delta\rVert^2_{\mathsf G}$,
which is \cref{eq:link-defect-local}. For the converse,
$\rho_x(g_{i_0}U_{i_0})^\top(R^x_{i_0})^{-\top}\widetilde{\mathsf G}_{i_0}(R^x_{i_0})^{-1}\rho_x(g_{i_0}U_{i_0})
=\rho_x(U_{i_0})^\top\widetilde{\mathsf G}_{i_0}\rho_x(U_{i_0})$.

(d) For $\mathbf B=\rho_x(\mathring\Theta^x_p)$ with $\mathbf B^\top\mathsf G\mathbf B=\mathsf G$,
$\Delta^\top\mathsf G\Delta=2\mathsf G-\mathbf B^\top\mathsf G-\mathsf G\mathbf B$ and
$\Tr(\mathsf G^{-1}\mathbf B^\top\mathsf G)=\Tr\mathbf B$, so
$2\mathfrak Q^x_p=2\dim\mathbb V-2\Tr\mathbf B$. The matrix $\mathsf G^{1/2}\mathbf B\mathsf G^{-1/2}$ is
orthogonal, so the eigenvalues of $\mathbf B$ have modulus one and
$\lvert\Tr\mathbf B\rvert\leq\dim\mathbb V$; $\Tr\mathsf U=\Tr\mathbf B$ because the two are similar.

(e) By linearity of $\mathrm d\rho_x$ and \cref{eq:link-generator-reversal},
$\sum_a\mathring\Omega_a=\mathrm d\rho_x(\mathbf C\mathring\omega^x)_p$. Expanding each $\exp(\mathring\Omega_a)$ in its
series and multiplying, $\rho_x(\mathring\Theta^x_p)-I_{\mathbb V}-\sum_a\mathring\Omega_a$ is the absolutely
convergent sum $\mathbf T_p$ of the ordered products
$\mathring\Omega_0^{l_0}\cdots\mathring\Omega_{r-1}^{l_{r-1}}/(l_0!\cdots l_{r-1}!)$ over the multi-indices with
$l_0+\cdots+l_{r-1}\geq2$, and by (a)
$\lVert\mathbf T_p\rVert_{\mathsf G}\leq\sum_{l_0+\cdots+l_{r-1}\geq2}\prod_a\lVert\mathring\Omega_a\rVert_{\mathsf G}^{l_a}/l_a!
=\exp(\mathring\varsigma_p)-1-\mathring\varsigma_p$. Since
$2\mathfrak Q^x_p=\lVert\sum_a\mathring\Omega_a+\mathbf T_p\rVert^2_{\mathsf G}$ and
$\lVert\sum_a\mathring\Omega_a\rVert_{\mathsf G}\leq\mathring\varsigma_p$, expanding the square gives
\cref{eq:link-defect-expansion}. From
$\exp(\mathring\varsigma)-1-\mathring\varsigma=\mathring\varsigma^2\sum_{k\geq0}\mathring\varsigma^k/(k+2)!\leq\tfrac12\mathring\varsigma^2\exp(\mathring\varsigma)$,
the right side of \cref{eq:link-defect-expansion} is at most
$\tfrac12\mathring\varsigma_p^3\exp(\mathring\varsigma_p)+\tfrac18\mathring\varsigma_p^4\exp(2\mathring\varsigma_p)$,
hence at most $(\tfrac12\exp(1)+\tfrac18\exp(2))\mathring\varsigma_p^3$ for $\mathring\varsigma_p\leq1$,
and $\tfrac12\exp(1)+\tfrac18\exp(2)<2.283$. With the pulled-back norm,
$\lVert\mathring\Omega_a\rVert_{\mathsf G}=\lVert\mathring\omega^x_{e_a}\rVert_{\mathsf G}\leq\lVert\mathring\omega^x\rVert$.
\hfill$\square$

\noindent\emph{Proof of \Cref{prop:link-generator-prior}.}
(a) By \Cref{prop:link-loop-defect}(a) each $\mathfrak Q^x_p$, and hence $\mathscr S$, is finite,
continuous, and nonnegative. For the mean, the series of the exponential and
\Cref{prop:link-loop-defect}(a) give
$\lVert\exp\mathbf H\rVert_{\mathsf G}\leq\lVert I_{\mathbb V}\rVert_{\mathsf G}+e^{\lVert\mathbf H\rVert_{\mathsf G}}-1\leq(\dim\mathbb V)^{1/2}e^{\lVert\mathbf H\rVert_{\mathsf G}}$,
so for a walk $p$ of length $r$, with $\mathring\varsigma_p$ as in \Cref{prop:link-loop-defect}(e),
\begin{equation*}
\mathfrak Q^x_p\leq\tfrac12\bigl((\dim\mathbb V)^{1/2}+(\dim\mathbb V)^{r/2}e^{\mathring\varsigma_p}\bigr)^2
\leq2(\dim\mathbb V)^re^{2\mathring\varsigma_p}.
\end{equation*}
The quantity $\mathring\varsigma_p$ is at most the operator norm of $\mathrm d\rho_x$ times the sum of
the norms of the generators on the pairs that $p$ traverses, each counted as often as it is traversed.
Under $\nu^{\mathscr T}$ these generators are independent Gaussian vectors, and the norm of a Gaussian
vector has finite exponential moments of every order, so $e^{2\mathring\varsigma_p}$ and therefore
$\mathscr S$ have finite means. Since $0<e^{-\mathscr S}\leq1$ and $\nu^{\mathscr T}$ is a probability law,
$0<Z_{\mathscr S}\leq1$, and convexity of $u\mapsto e^{-u}$ gives
$Z_{\mathscr S}=\E_{\nu^{\mathscr T}}e^{-\mathscr S}\geq\exp(-\E_{\nu^{\mathscr T}}\mathscr S)$ by Jensen's
inequality. When $\rho_x(G)$ consists of isometries, \Cref{prop:link-loop-defect}(d) gives
$\mathscr S\leq2\dim\mathbb V\sum_p\mathrm c_p$. The density is $Z_{\mathscr S}^{-1}e^{-\mathscr S}$
times the product of the Gaussian densities, which are positive and continuous because every
covariance is positive definite. The hypotheses of \Cref{prop:elbo-random-link-vfe} on base law, action,
and normalizer are that $\nu^{\mathscr T}$ is a probability law on a standard Borel space, that
$\mathscr S$ is measurable with values in $(-\infty,+\infty]$, and that $0<Z_{\mathscr S}<\infty$. The
defect depends on the generators only through the links (\Cref{prop:link-loop-defect}(a)), and
\cref{eq:link-holonomy-chart} gives $\mathring\Theta^x_p=U_{i_0}^{-1}\mathrm{Hol}^x(p)U_{i_0}$, so
$\widetilde{\mathscr S}$ is a continuous function on $G^{E_+}$ with
$\mathscr S=\widetilde{\mathscr S}\circ\exp_\sigma$, and the change of variables for image measures gives
$\int e^{-\widetilde{\mathscr S}}d(\exp_\sigma)_\#\nu^{\mathscr T}=Z_{\mathscr S}$.

(b) The generators, anchors, covariances, inner product, walks, and the matrix at $\sigma$ are unchanged
by passive re-expression, and the links move by the vertex law (\Cref{prop:link-generator-gauge}(a)).

(c) By \cref{eq:link-exp-naturality},
$\mathrm d\rho_x(\operatorname{Ad}_{g_0^{-1}}\mathring\omega)=\mathbf R^{-1}\mathrm d\rho_x(\mathring\omega)\mathbf R$,
so polarizing \cref{eq:link-defect-conjugation} with $\mathsf G'=\mathbf R^\top\mathsf G^{\mathrm{fib}}\mathbf R$
shows that $\operatorname{Ad}_{g_0^{-1}}$ maps the pullback for $\mathsf G^{\mathrm{fib}}$ isometrically onto
the pullback for $\mathbf R^\top\mathsf G^{\mathrm{fib}}\mathbf R$; in general $\operatorname{Ad}_{g_0^{-1}}$
is an isometry from $\langle\cdot,\cdot\rangle_{\mathsf G}$ onto
$\langle\operatorname{Ad}_{g_0}\cdot,\operatorname{Ad}_{g_0}\cdot\rangle_{\mathsf G}$ by construction. A linear
isometry between inner product spaces carries the normalized Lebesgue measure to the normalized Lebesgue
measure and $\mathcal N(\mathring\alpha,\mathring\Sigma)$ to the Gaussian law with mean
$\operatorname{Ad}_{g_0^{-1}}\mathring\alpha$ and covariance operator
$\operatorname{Ad}_{g_0^{-1}}\mathring\Sigma\operatorname{Ad}_{g_0}$: in the new inner product
$\langle\operatorname{Ad}_{g_0}\cdot,\operatorname{Ad}_{g_0}\cdot\rangle_{\mathsf G}$, the pairing of
$\mathring\beta\in\mathfrak g$ with $\operatorname{Ad}_{g_0^{-1}}\mathring\omega$ is
$\langle\operatorname{Ad}_{g_0}\mathring\beta,\mathring\omega\rangle_{\mathsf G}$, whose variance is
$\langle\operatorname{Ad}_{g_0}\mathring\beta,\mathring\Sigma\operatorname{Ad}_{g_0}\mathring\beta\rangle_{\mathsf G}$. By \Cref{prop:link-loop-defect}(b) the defect at $\sigma'$ with
$\mathbf R^\top\mathsf G^{\mathrm{fib}}\mathbf R$ of $\operatorname{Ad}_{g_0^{-1}}\mathring\omega^x$ equals
$\mathfrak Q^x_p(\mathring\omega^x)$, so the transformed action composed with
$\operatorname{Ad}_{g_0^{-1}}$ is $\mathscr S$ and the normalizer is unchanged. Finally
$\exp_{\sigma'}\circ\operatorname{Ad}_{g_0^{-1}}=\exp_\sigma$ by \Cref{prop:link-generator-gauge}(b).

(d) Under the hypotheses, the data of (c) are the anchors, covariances, inner product, walks, and
coefficients of the prior together with the fiber matrix
$\mathbf R^\top\mathsf G^{\mathrm{fib}}\mathbf R=\mathrm c_{\mathrm o}^2\mathsf G^{\mathrm{fib}}$ for a real
$\mathrm c_{\mathrm o}\neq0$, and the trace form \cref{eq:link-trace-form}, hence the defect,
is unchanged when its matrix is multiplied by a positive scalar (\Cref{lem:link-trace-form-scale}). For
the pulled-back inner product,
$\lVert\mathrm d\rho_x(\operatorname{Ad}_{g_0^{-1}}\mathring\omega)\rVert_{\mathsf G}
=\lVert\mathrm d\rho_x\mathring\omega\rVert_{\mathbf R^{-\top}\mathsf G^{\mathrm{fib}}\mathbf R^{-1}}$ by
\cref{eq:link-defect-conjugation}, which equals $\lVert\mathrm d\rho_x\mathring\omega\rVert_{\mathsf G}$ when
$\mathbf R^{-\top}\mathsf G^{\mathrm{fib}}\mathbf R^{-1}=\mathrm c_{\mathrm o}^{-2}\mathsf G^{\mathrm{fib}}$; when
$\mathrm d\rho_x(\mathfrak g)=\mathfrak{gl}(\mathbb V)$, equality for every $\mathring\omega$ forces
$\mathbf R^{-\top}\mathsf G^{\mathrm{fib}}\mathbf R^{-1}$ to be a positive multiple of
$\mathsf G^{\mathrm{fib}}$ by \Cref{lem:link-trace-form-scale}. \hfill$\square$

\noindent\emph{Proof of \Cref{prop:link-translation-instance}.}
(a) Block multiplication gives $\rho(\mathring\omega)\rho(\mathring\omega')=\rho(\mathring\omega+\mathring\omega')$,
and $\rho$ is injective and smooth. The one-parameter subgroups of $(\R^K,+)$ are
$t\mapsto t\mathring\omega$, so $\exp$ is the identity, and $\mathrm d\rho(\mathring\omega)^2=0$ gives
$\exp(\mathrm d\rho\mathring\omega)=I_{K+1}+\mathrm d\rho(\mathring\omega)=\rho(\mathring\omega)$. The
twist \cref{eq:geo-edge-twist} reads $\mathring\omega_e=-U_i+\Theta_e+U_j$ additively, and along a closed
walk the frame terms $U_{i_{a+1}}-U_{i_a}$ sum to zero.

(b) With $\mathbf b$ the last standard basis vector of $\R^{K+1}$ and $\mathring\xi=(\mathbf C\mathring\omega)_p$, the
matrix $I_{K+1}-\rho(\mathring\xi)=-\mathrm d\rho(\mathring\xi)=-(\mathring\xi,0)\mathbf b^\top$ has rank at most one, and the rank-one
formula in the proof of \Cref{lem:link-trace-form-scale} gives
$\lVert(\mathring\xi,0)\mathbf b^\top\rVert^2_{\mathsf G}=(\mathring\xi,0)^\top\mathsf G^{\mathrm{fib}}(\mathring\xi,0)\cdot\mathbf b^\top(\mathsf G^{\mathrm{fib}})^{-1}\mathbf b$,
which is both $2\mathfrak Q_p$ and $\lVert\mathring\xi\rVert^2_{\mathsf G}$.

(c) The configuration priors are stated in \Cref{def:rg-formation-cycle-experiment,def:rg-gated-translation-family},
and a standard normal vector has independent coordinates. In generator coordinates
$\mathring\omega_e=\omega_e-U_i+U_j$ is normal with mean $U_j-U_i$ and unit variance, and around a closed
walk the means telescope to zero. The edge $e_i$ points from $i$ to $i+1$, so $e_a:i_{a+1}\to i_a$ holds
along $(e_3,e_2,e_1,e_0)$ with the vertices $0,3,2,1,0$, and part (b) gives the defect. The residual sums
telescope the means around the cycle.

(d) The normalized loop covector of $q$ has the entries $1/\sqrt3$ on the links of $q$ and zero
elsewhere (\Cref{lem:rg-formation-carrier}), so $h_q=\mathfrak h_q/\sqrt3$, and part (b) gives
$\mathfrak Q_q=\tfrac12\mathfrak h_q^2$, hence $\mathrm J_1h_q^2=\tfrac23\mathrm J_1\mathfrak Q_q$. By
\cref{eq:rg-formation-physical-prior,eq:rg-formation-orthogonal-split},
$\Omega_0(\vartheta)=\Omega^\perp_0(\vartheta)+\vartheta\mathrm J_1P_{\mathcal H}$,
and $\mathbf x^\top P_{\mathcal H}\mathbf x=\sum_qh_q^2$ because the normalized loop covectors are
orthonormal, which gives \cref{eq:link-triangle-chain-prior}. By \Cref{lem:rg-formation-carrier} the rows
of $B$ and the loop covectors are orthogonal and the configuration space is the orthogonal sum
$\mathcal U\oplus\mathcal H\oplus\mathcal K$. The matrices $B^\top\operatorname{diag}(\mathrm J_\ell I)B$,
$P_{\mathcal K}$, and $R_{\mathrm{ev}}=\mathfrak D^\top\mathfrak D/(1-\sigma^2_{\mathrm{ev}})$, whose rows
$\mathfrak d_i$ are rows of $B$, annihilate $\mathcal H$ and map into $\mathcal U\oplus\mathcal K$, so
$\Omega^\perp_0(\vartheta)$ annihilates $\mathcal H$ and agrees with $\Omega_0(\vartheta)$ on
$\mathcal U\oplus\mathcal K$, where it is positive definite by
\Cref{prop:rg-formation-normalized-experiment}. Hence $\Omega_0(\vartheta)$ is block diagonal for
$\mathcal H\oplus(\mathcal U\oplus\mathcal K)$ with the block $\vartheta\mathrm J_1$ on $\mathcal H$, and
the Gaussian prior with precision $\Omega_0(\vartheta)/\mathrm T$ makes $P_{\mathcal H}\mathbf x$
independent of the other component, with covariance $(\mathrm T/(\vartheta\mathrm J_1))P_{\mathcal H}$;
since $\mathfrak h_q=\sqrt3h_q$ and the loop covectors are orthonormal, the $\mathfrak h_q$ are independent
with variance $3\mathrm T/(\vartheta\mathrm J_1)$. In \cref{eq:rg-formation-partition-integrals} the
matrix of $\mathrm E_\vartheta$ differs from $\vartheta H$ by a multiple of $P_{\mathcal K}$, its linear
term is along $z_\star\in\mathcal K$, and $\mathrm E^{\mathrm{phys}}_\vartheta$ differs from
$\mathrm E_\vartheta$ by a multiple of $\mathbf x^\top R_{\mathrm{ev}}\mathbf x$, so the same block
decomposition gives the energy statement. Finally, $\Omega^\perp_0(\vartheta)$ has the nonzero kernel
$\mathcal H$, so $e^{-\mathbf x^\top\Omega^\perp_0(\vartheta)\mathbf x/(2\mathrm T)}$ is constant along
$\mathcal H$ and not integrable, and it depends on means and links through $B$. \hfill$\square$

\noindent\emph{Proof of \Cref{prop:link-latent-normalization}.}
(a) The prior is a probability law by \cref{prop:link-generator-prior}. Composing the measurable map
$\mathring\omega\mapsto\Theta(\mathring\omega)$ with the probability kernel
$\Theta\mapsto\mathbb P^{\mathsf Y}_{\theta,V}(\cdot\given X,\Theta)$ of \cref{prop:gen-exact-normalization}
and with the record kernels gives probability kernels in $\mathring\omega$. Placing the generator family
before every agent in a topological order of $\Gamma$, as a parent of every factor that reads a link, makes
\eqref{eq:link-latent-joint} the iterated integral \eqref{eq:gen-iterated-integral} of a directed latent law
on $\mathscr T\times\mathsf Y$ with records conditionally independent given $(\mathring\omega,y)$, and
\cref{prop:gen-exact-normalization} gives unit mass on every finite record hypergraph. (b) The function
$L_o$ is jointly measurable, so Tonelli's theorem gives the measurability of $p(o\given\cdot)$ and, after
integrating $y$ before $\mathring\omega$, \eqref{eq:link-latent-evidence}. Since $\mathscr S$ is finite, the
prior has a strictly positive density with respect to Lebesgue measure on $\mathscr T$, so the integral of
the nonnegative function $p(o\given\cdot)$ is positive exactly when that function is positive on a set of
positive Lebesgue measure. Bounded densities give $L_o\leq\mathrm c_{\mathfrak A}$, and integration against probability
kernels gives both bounds; a binary success record has densities $e^{-E}\leq1$ and $1-e^{-E}\leq1$ because
$E\geq0$. (c) By \eqref{eq:obs-global-posterior} at $(X,\Theta(\mathring\omega))$, the measure
$L_o(\mathring\omega,y)\mathbb P^{\mathsf Y}(dy\given\mathring\omega)\mathbb P^{\mathscr T}_\theta(d\mathring\omega\given X)$
equals $p(o\given\mathring\omega)\boldsymbol\Pi_{\mathring\omega}(dy)\mathbb P^{\mathscr T}_\theta(d\mathring\omega\given X)$,
and dividing by the evidence gives the disintegration. \Cref{lem:prob-kl-chain-rule}(a) gives
$\KL(Q\Vert\boldsymbol\Pi_{\theta,V,o,X})=\KL(R\Vert\boldsymbol\Pi^{\mathscr T})+\int\KL(Q_{\mathring\omega}\Vert\boldsymbol\Pi_{\mathring\omega})R(d\mathring\omega)$
in $[0,+\infty]$ with a measurable integrand, and the evidence lies in $(0,\mathrm c_{\mathfrak A}]$ by (b). If
$\KL(R\Vert\mathbb P^{\mathscr T}_\theta)<\infty$, then
$R\ll\mathbb P^{\mathscr T}_\theta$, and since
$d\boldsymbol\Pi^{\mathscr T}/d\mathbb P^{\mathscr T}_\theta=p(o\given\cdot)/p_{\theta,V}(o\given X)$ is positive
and finite, $R\ll\boldsymbol\Pi^{\mathscr T}$ with
$\log(dR/d\boldsymbol\Pi^{\mathscr T})=\log(dR/d\mathbb P^{\mathscr T}_\theta)-\log p(o\given\mathring\omega)+\log p_{\theta,V}(o\given X)$.
The first term has an $R$-integrable negative part (\cref{def:prob-log-density-ratio-convention}) and the
second is bounded below by $-\log\mathrm c_{\mathfrak A}$, so the integral splits, and adding the conditional term and
$-\log p_{\theta,V}(o\given X)$ gives \eqref{eq:link-decomposition}, because
$\Fenergy_V[Q_{\mathring\omega};X,\Theta(\mathring\omega),o]=-\log p(o\given\mathring\omega)+\KL(Q_{\mathring\omega}\Vert\boldsymbol\Pi_{\mathring\omega})$.
If $\KL(R\Vert\mathbb P^{\mathscr T}_\theta)=+\infty$, the bound
$\boldsymbol\Pi^{\mathscr T}\leq(\mathrm c_{\mathfrak A}/p_{\theta,V}(o\given X))\mathbb P^{\mathscr T}_\theta$ gives
$dR/d\boldsymbol\Pi^{\mathscr T}\geq(p_{\theta,V}(o\given X)/\mathrm c_{\mathfrak A})dR/d\mathbb P^{\mathscr T}_\theta$ whenever
$R\ll\boldsymbol\Pi^{\mathscr T}$, so $\KL(R\Vert\boldsymbol\Pi^{\mathscr T})=+\infty$ and both sides of
\eqref{eq:link-decomposition} are $+\infty$. The last sentence is
$\log(dR/d\mathbb P^{\mathscr T}_\theta)=\log(dR/d\nu^{\mathscr T})+\mathscr S+\log Z_{\mathscr S}$ integrated
under the stated integrability. \hfill$\square$

\noindent\emph{Proof of \Cref{lem:link-gaussian-gradients}.}
All expectations below are finite by \cref{lem:link-growth}(a). First, for a continuously differentiable
$\mathbf t$ with $\lvert\mathbf t\rvert+\lVert\nabla\mathbf t\rVert\leq\mathrm c_1e^{\mathrm c_2\lVert\mathring\omega\rVert}$,
\begin{equation}
\E_{\mathcal N(\mu,\Sigma)}\bigl[(\mathring\omega-\mu)\mathbf t(\mathring\omega)\bigr]=\Sigma\E_{\mathcal N(\mu,\Sigma)}\nabla\mathbf t.
\label{eq:link-stein}
\end{equation}
Indeed, for a standard Gaussian vector $\mathbf b$ in an orthonormal basis and a continuously
differentiable $\mathbf t'$ of exponential growth with its gradient, integration by parts in the $l$-th
coordinate against the standard normal density, whose derivative is minus the coordinate times the density,
has vanishing boundary terms and gives $\E[\mathbf b_l\mathbf t'(\mathbf b)]=\E[\partial_l\mathbf t'(\mathbf b)]$
after Fubini's theorem; with $\mathbf t'(\mathbf b)=\mathbf t(\mu+\Sigma^{1/2}\mathbf b)$, whose gradient is
$\Sigma^{1/2}\nabla\mathbf t$, this is \eqref{eq:link-stein}. The mean derivative follows by differentiating
$\E\mathbf h(\mu+\Sigma^{1/2}\mathbf b)$ under the expectation, with domination by
$\mathrm c_1e^{\mathrm c_2(\lVert\mu\rVert+\lVert\mathring\eta\rVert+\lVert\Sigma\rVert^{1/2}\lVert\mathbf b\rVert)}$ along
$\mu+t\mathring\eta$, $\lvert t\rvert\leq1$, for each direction $\mathring\eta$. For the covariance, $\Sigma_t=\Sigma+t\mathring\Delta$ stays between half the
smallest eigenvalue of $\Sigma$ and $\lVert\Sigma\rVert$ plus that amount for small $\lvert t\rvert$, so the
$t$-derivative of the Gaussian density,
$\frac12[\langle\mathring\omega-\mu,\Sigma_t^{-1}\mathring\Delta\Sigma_t^{-1}(\mathring\omega-\mu)\rangle-\Tr(\Sigma_t^{-1}\mathring\Delta)]$
times the density, is dominated uniformly by a Gaussian times a quadratic, and against
$\lvert\mathbf h\rvert\leq\mathrm c_1e^{\mathrm c_2\lVert\mathring\omega\rVert}$ the derivative passes under the integral:
$\frac{d}{dt}\E\mathbf h|_{t=0}=\frac12\E[\mathbf h(\mathring\omega)\langle\mathring\omega-\mu,\Sigma^{-1}\mathring\Delta\Sigma^{-1}(\mathring\omega-\mu)\rangle]-\frac12\Tr(\Sigma^{-1}\mathring\Delta)\E\mathbf h$.
Applying \eqref{eq:link-stein} to $\mathbf t=\mathbf h\langle\mathring\xi,\mathring\omega-\mu\rangle$ for a given
vector $\mathring\xi$, and then to the directional derivatives of $\mathbf h$, gives
$\E[\mathbf h(\mathring\omega-\mu)\otimes(\mathring\omega-\mu)]=\Sigma(\E\nabla^2\mathbf h)\Sigma+\Sigma\E\mathbf h$,
so the first term equals $\frac12\Tr(\mathring\Delta\E\nabla^2\mathbf h)+\frac12\Tr(\Sigma^{-1}\mathring\Delta)\E\mathbf h$, which is
\eqref{eq:link-bonnet-price}. Both derivatives are continuous in $(\mu,\Sigma)$ by dominated convergence in
the representation $\mathring\omega=\mu+\Sigma^{1/2}\mathbf b$, since $\Sigma\mapsto\Sigma^{1/2}$ is continuous on
the cone. \hfill$\square$

\noindent\emph{Proof of \Cref{prop:link-field-equation}.}
(a) The action satisfies \eqref{eq:link-growth} by \cref{lem:link-growth}(c). On each ball
$\lVert\mathring\omega\rVert\leq\mathring r$, \eqref{eq:link-matter-domination} dominates $\mathbf H_o$ and its first two
$\mathring\omega$-derivatives by an integrable function of $y$, so differentiation passes under
$\E_{Q^{\mathsf Y}}$ and $\E_{Q^{\mathsf Y}}\mathbf H_o$ is twice continuously differentiable with derivatives bounded
by $\mathrm c_1e^{\mathrm c_2\mathring r}$ on that ball. At every $\mathring\omega$ the held-link model at
$\Theta(\mathring\omega)$ satisfies \cref{hyp:elbo-evidence-domain} for $Q^{\mathsf Y}$: its posterior has the
density proportional to $e^{-\mathbf H_o(\mathring\omega,\cdot)}$, which is positive wherever $\mathbf H_o$ is
finite, hence $Q^{\mathsf Y}$-almost everywhere by \eqref{eq:link-matter-domination}, and both halves of
\eqref{eq:elbo-log-integrability}--\eqref{eq:prob-integrability} are finite. So
$\Fenergy_V[Q^{\mathsf Y};X,\Theta(\mathring\omega),o]=\E_{Q^{\mathsf Y}}[\log q^{\mathsf Y}+\mathbf H_o(\mathring\omega,Y)]$ by
\cref{thm:elbo-exact-identity}, and \cref{prop:link-latent-normalization}(c) with
\cref{prop:link-gaussian-recognition-vfe}(a,b) gives \eqref{eq:link-field-free-energy}, finite because
$\E_R\lvert\mathbf U\rvert<\infty$ by \cref{lem:link-growth}(a). The closed form \eqref{eq:link-gaussian-kl} is
smooth, and $\E_R\mathbf U$ is continuously differentiable in the mean and the covariance
$\mathring\Lambda^{-1}$ by \cref{lem:link-gaussian-gradients}; the inversion
$\mathring\Lambda\mapsto\mathring\Lambda^{-1}$ is a diffeomorphism of the cone.

(b) By \eqref{eq:link-gaussian-kl} and \cref{lem:link-gaussian-gradients}, the derivative of
\eqref{eq:link-field-free-energy} in the mean along $\mathring\eta$ is
$\langle\mathring\Sigma^{-1}(\mathring\mu-\mathring\alpha)+\E_R\nabla\mathbf U,\mathring\eta\rangle$, and its derivative in the
covariance along a self-adjoint $\mathring\Delta$ is
$\frac12\Tr\bigl(\mathring\Delta(\mathring\Sigma^{-1}-\mathring\Lambda+\E_R\nabla^2\mathbf U)\bigr)$. Stationarity in
$(\mathring\mu,\mathring\Lambda)$ is stationarity in $(\mathring\mu,\mathring\Lambda^{-1})$, and a self-adjoint
operator orthogonal to every self-adjoint $\mathring\Delta$ vanishes; differentiation under $\E_{Q^{\mathsf Y}}$
gives $\E_R\nabla\mathbf U=\E_R\nabla\mathscr S+\E_{R\otimes Q^{\mathsf Y}}\nabla_{\mathring\omega}\mathbf H_o$ and the same
for Hessians. In the mean-field case the admissible $\mathring\Delta$ are block diagonal, and $\mathring\Sigma$ and
$\mathring\Lambda$ are block diagonal, so only the diagonal blocks are constrained.

(c) For given $y$ and given remaining generators, \cref{lem:link-exp-derivative}(c) applied to
$\mathbf H_o$ as a function of the matrix $\rho_x(\mathring\Theta^x_e)$ gives
$\nabla_{(x,e)}\mathbf H_o=(\mathrm d\rho_x)^\top\mathbf L\bigl((\mathrm d\rho_x\mathring\omega^x_e)^\top,\partial\mathbf H_o/\partial\rho_x(\mathring\Theta^x_e)\bigr)$;
the map $\mathbf S\mapsto(\mathrm d\rho_x)^\top\mathbf L(\mathbf B,\mathbf S)$ is linear and continuous, so it
commutes with $\E_{Q^{\mathsf Y}}$, and the $(x,e)$ component of \eqref{eq:link-field-mean} becomes
\eqref{eq:link-field-current}.

(d) Since $\E_R(\mathring\omega-\mathring\alpha)=\mathring\mu-\mathring\alpha$, \eqref{eq:link-field-mean} is
$\E_R\nabla\mathbf V=0$. With $\mathring\omega=\mathring\mu+\mathring\Lambda^{-1/2}\mathbf b$, $\nabla\mathbf V$ is uniformly
continuous on bounded sets and dominated by
$\mathrm c(1+\lVert\mathring\mu\rVert+\lVert\mathbf b\rVert)e^{\mathrm c(\lVert\mathring\mu\rVert+\lVert\mathbf b\rVert)}$ for
$\lVert\mathring\Lambda^{-1}\rVert\leq1$; splitting the expectation at $\lVert\mathbf b\rVert=\mathrm c'$ bounds the tail
uniformly for $\mathring\mu$ in a compact set, and the modulus of continuity controls the rest, which gives
the uniform convergence. A limit point then satisfies
$\lVert\nabla\mathbf V(\mathring\mu_*)\rVert\leq\lVert\nabla\mathbf V(\mathring\mu_*)-\nabla\mathbf V(\mathring\mu)\rVert
+\lVert\nabla\mathbf V(\mathring\mu)-\E_R\nabla\mathbf V\rVert\to0$ along the solutions. For the curve, the map
$(\mathring\mu,t)\mapsto\E\nabla\mathbf V(\mathring\mu+t\mathring\Lambda_0^{-1/2}\mathbf b)$ is continuously
differentiable on $\mathscr T\times\R$ by dominated differentiation, equals $\nabla\mathbf V(\mathring\mu)$ at
$t=0$, and has $\mathring\mu$-derivative $\nabla^2\mathbf V(\mathring\mu_*)$ at $(\mathring\mu_*,0)$, so the implicit
function theorem gives $t\mapsto\mathring\mu(t)$. For the last claim, every law $R$ with
$\KL(R\Vert\nu^{\mathscr T})=+\infty$ has $\KL(R\Vert\mathbb P^{\mathscr T}_\theta)=+\infty$, because
$\mathbb P^{\mathscr T}_\theta\leq Z_{\mathscr S}^{-1}\nu^{\mathscr T}$, and hence infinite free energy by
\eqref{eq:link-decomposition}; for the others, \eqref{eq:link-decomposition} and the split used in (a) give
\eqref{eq:link-field-free-energy} in $(-\infty,+\infty]$, because the matter integrand is bounded below by
$-\log\mathrm c_{\mathfrak A}$ and $\mathscr S\geq0$, and with
$0<\int e^{-\mathbf U}d\nu^{\mathscr T}<\infty$,
$\KL(R\Vert\nu^{\mathscr T})+\E_R\mathbf U=\KL(R\Vert R^\star)-\log\int e^{-\mathbf U}d\nu^{\mathscr T}$, which is
minimal exactly at $R=R^\star$. The density of $\nu^{\mathscr T}$ is proportional to
$e^{-\frac12\langle\mathring\omega-\mathring\alpha,\mathring\Sigma^{-1}(\mathring\omega-\mathring\alpha)\rangle}$.

(e) A singleton is null for $\nu^{\mathscr T}$, which has a Lebesgue density, so
\eqref{eq:prob-direct-null-set} of \cref{prop:prob-density-absolute-continuity} gives the infinite
divergence, and the free energy is infinite as in (d). At a given $\mathring\mu$, \eqref{eq:link-gaussian-kl}
is at least $\frac12(\log\det\mathring\Lambda-\log\det\mathring\Sigma^{-1}-\dim\mathscr T)$, which tends to
$+\infty$ as $\mathring\Lambda^{-1}\to0$, while $\E_R\mathbf U\to\mathbf U(\mathring\mu)$ by the argument of (d)
applied to $\mathbf U$. \hfill$\square$

\noindent\emph{Proof of \Cref{prop:link-law-level-energy}.}
(a) Pushforward by the bimeasurable map $\rho_x(h)$ preserves relative entropy
(\cref{lem:prob-kl-chain-rule}(e)), so
$D_e(h\Theta h')=\KL(q_i\Vert\widehat\rho_x(h)\widehat\rho_x(\Theta)q_j)=\KL(\widehat\rho_x(h)^{-1}q_i\Vert\widehat\rho_x(\Theta)q_j)=D_e(\Theta)$.
The maps $\Theta\mapsto h\Theta h'$ are diffeomorphisms of $G$ that preserve $D_e$, so they preserve its
minimizers and, where it is continuously differentiable, its stationary points. If $0\neq\mathring\xi\in\mathfrak g_{q_i}$,
the curve $t\mapsto\exp(t\mathring\xi)\Theta$ lies in $G_{q_i}\Theta$ and is not constant near $t=0$, since the exponential
map is injective near $0$; the case of $\mathfrak g_{q_j}$ uses $\Theta\exp(t\mathring\xi)$. If the exponential map is a
local diffeomorphism at the generator, so is the chart map $\mathring\omega^x_e\mapsto U_i\exp(\mathring\omega^x_e)U_j^{-1}$,
and its local inverse carries such a curve to a non-constant curve of minimizers or stationary points of
the chart energy. (b) The pushforward of $\mathcal N(\mu,\Sigma)$ by $g$ is $\mathcal N(g\mu,g\Sigma g^\top)$.
Writing $g=\Sigma^{1/2}O\Sigma^{-1/2}$, the condition $g\Sigma g^\top=\Sigma$ is $OO^\top=I_K$, $\det g=\det O>0$
gives $O\in\mathrm{SO}(K)$, and $g\mu=\mu$ is $O\Sigma^{-1/2}\mu=\Sigma^{-1/2}\mu$; the stabilizer in $\mathrm{SO}(K)$ of a
nonzero vector is a copy of $\mathrm{SO}(K-1)$. Invariance of the Mahalanobis norm is
$(g\mu)^\top(g\Sigma g^\top)^{-1}g\mu=\mu^\top\Sigma^{-1}\mu$, which proves necessity. For sufficiency, $\mathrm{SO}(K)$ acts
transitively on spheres of $\R^K$ when $K\geq2$, so some $O$ carries $\Sigma_j^{-1/2}\mu_j$ to
$\Sigma_i^{-1/2}\mu_i$, and $\Theta_0=\Sigma_i^{1/2}O\Sigma_j^{-1/2}$ transports $q_j$ to $q_i$. The zero set of
$D_e$ is $\{\Theta:\widehat\rho_x(\Theta)q_j=q_i\}=\Theta_0G_{q_j}$, and $\Theta_0^{-1}G_{q_i}\Theta_0=G_{q_j}$. The
dimension count and (a) give the last sentence. (c) Under \cref{hyp:pb-regular-models} relative entropy is
$C^3$ near the diagonal (\cref{prop:pb-kl-divergence-jets}) and $\Theta\mapsto\widehat\rho_x(\Theta)q_j$ is smooth,
so $D_e$ is $C^3$ near $\Theta_0$; it is nonnegative and vanishes at $\Theta_0$, so its differential vanishes
there and its Hessian is a well-defined form. Along $t\mapsto\exp(t\mathring\xi)\Theta_0$,
$D_e=\KL(q_i\Vert\widehat\rho_x(\exp t\mathring\xi)q_i)=\frac{t^2}2\mathfrak k_{q_i}(\mathring\xi,\mathring\xi)+O(t^3)$ by
\eqref{eq:pb-orbit-kl-expansion}, so the Hessian on the velocity $\mathring\xi\Theta_0$ is $\mathfrak k_{q_i}(\mathring\xi,\mathring\xi)$; every
tangent vector at $\Theta_0$ has this form, the radical of $\mathfrak k_{q_i}$ is $\mathfrak g_{q_i}$
(\cref{cor:pb-orbit-fisher-form}), and $\mathfrak g_{q_i}\Theta_0$ is the tangent space of $G_{q_i}\Theta_0$. For
Gaussian laws the expansion can be computed directly: with
$\mu_t=e^{t\mathring\xi}\mu=\mu+t\mathring\xi\mu+O(t^2)$, $\Sigma_t=e^{t\mathring\xi}\Sigma e^{t\mathring\xi^\top}$, and
$\Sigma^{-1/2}\Sigma_t\Sigma^{-1/2}=I_K+\mathbf T_t$, where $\mathbf T_t=t\mathbf B+O(t^2)$ is self-adjoint and
$\mathbf B=\Sigma^{-1/2}\mathring\xi\Sigma^{1/2}+(\Sigma^{-1/2}\mathring\xi\Sigma^{1/2})^\top$, the Gaussian divergence
$\KL(\mathcal N(\mu,\Sigma)\Vert\mathcal N(\mu_t,\Sigma_t))$ is
$\frac12[\Tr((I_K+\mathbf T_t)^{-1})-K+\log\det(I_K+\mathbf T_t)]+\frac12(\mu_t-\mu)^\top\Sigma_t^{-1}(\mu_t-\mu)$.
The power series of $(I_K+\mathbf T_t)^{-1}$ and of $\log\det(I_K+\mathbf T_t)=\Tr\log(I_K+\mathbf T_t)$, valid
for $\lVert\mathbf T_t\rVert<1$, make the first term $\frac14\Tr(\mathbf T_t^2)+O(t^3)=\frac{t^2}4\Tr(\mathbf B^2)+O(t^3)$,
since its terms of first order cancel, and the second term is
$\frac{t^2}2\lvert\Sigma^{-1/2}\mathring\xi\mu\rvert^2+O(t^3)$, which gives \eqref{eq:link-gaussian-orbit-form}. The form
vanishes exactly when $\mathring\xi\mu=0$ and $\Sigma^{-1/2}\mathring\xi\Sigma^{1/2}$ is skew, the Lie algebra of $G_q$ in (b).
\hfill$\square$

\noindent\emph{Proof of \Cref{cor:link-prior-necessary}.}
(a)(i) With the second moment $\Sigma_j+\mu_j\mu_j^\top\succ0$ and $\E[k_ik_j^\top]=\mu_i\mu_j^\top$, the gradient is
$\nabla\mathbf l(\mathbf M)=-W_a\bigl(\mu_i\mu_j^\top-\mathbf M(\Sigma_j+\mu_j\mu_j^\top)\bigr)$, which vanishes exactly at
$\mathbf M^\star$, and
\begin{equation}
\mathbf l(\mathbf M)-\mathbf l(\mathbf M^\star)
=\tfrac12\bigl\lVert W_a^{1/2}(\mathbf M-\mathbf M^\star)(\Sigma_j+\mu_j\mu_j^\top)^{1/2}\bigr\rVert_{\mathrm F}^2,
\label{eq:link-strong-convexity}
\end{equation}
so $\mathbf l$ is strictly convex; the second form of \eqref{eq:link-rank-one} is the Sherman--Morrison
identity, and the rank bound is that of $\mu_i\mu_j^\top$. With a cross covariance the same computation uses
$\E[k_ik_j^\top]=\mu_i\mu_j^\top+\Sigma_{ij}$. (ii) A rank at most one matrix is not invertible for $K\geq2$. Every
singular matrix is a limit of matrices of positive determinant: in a singular value decomposition, replace
each zero singular value by a small positive number, and change the sign of one of them if the determinant
would be negative. Continuity gives the infimum, uniqueness of the minimizer excludes attainment, and
\eqref{eq:link-strong-convexity} forces minimizing sequences to converge. (iii) By
\cref{lem:link-exp-derivative}(c) the gradient of $\mathring\omega\mapsto\mathbf l(e^{\mathring\omega})$ for the trace
pairing is $\mathbf L(\mathring\omega^\top,\nabla\mathbf l(e^{\mathring\omega}))$, and $\nabla\mathbf l(e^{\mathring\omega})\neq0$ by (ii), so
stationarity forces $\mathbf L(\mathring\omega^\top,\cdot)$, and with it its adjoint $\mathbf L(\mathring\omega,\cdot)$, to be
singular. In the example, with
$\mathbf R_{\mathring\phi}=\exp\bigl(\mathring\phi\begin{psmallmatrix}0&-1\\1&0\end{psmallmatrix}\bigr)$ the rotation by the
angle $\mathring\phi$, $e^{\mathring\omega}=\frac13\mathbf R_{2\pi}=\frac13I_2$ and
$\nabla\mathbf l(\frac13I_2)=-\bigl(\mu_i\mu_j^\top-\frac13(I_2+\mu_j\mu_j^\top)\bigr)=\frac13\operatorname{diag}(-1,1)$;
this matrix composed with $\mathbf R_{\mathring\phi}$ equals $\mathbf R_{-\mathring\phi}$ composed with it, and since
$e^{t\mathring\omega^\top}=3^{-t}\mathbf R_{-2\pi t}$ and $e^{(1-t)\mathring\omega^\top}=3^{t-1}\mathbf R_{2\pi t}$,
$\mathbf L\bigl(\mathring\omega^\top,\operatorname{diag}(-1,1)\bigr)=\frac13\int_0^1\mathbf R_{-4\pi t}dt\operatorname{diag}(-1,1)=0$.
(iv) If $\mu_j^\top\Sigma_j^{-1}\mu_i>0$, then $\mathbf M^\star$ has the simple eigenvalue
$\operatorname{tr}\mathbf M^\star>0$ and the eigenvalue $0$ with a $(K-1)$-dimensional eigenspace, so
$\mathbf M^\star=\mathbf S\operatorname{diag}(\operatorname{tr}\mathbf M^\star,0,\ldots,0)\mathbf S^{-1}$ for a real invertible
$\mathbf S$, and for $\mathrm c_0>0$ the chart links
$\mathbf S\operatorname{diag}(\operatorname{tr}\mathbf M^\star,\mathrm c_0,\ldots,\mathrm c_0)\mathbf S^{-1}
=\exp\bigl(\mathbf S\operatorname{diag}(\log\operatorname{tr}\mathbf M^\star,\log\mathrm c_0,\ldots,\log\mathrm c_0)\mathbf S^{-1}\bigr)$
converge to it as $\mathrm c_0\to0$. If $\mu_j^\top\Sigma_j^{-1}\mu_i=0$, then $(\mathbf M^\star)^2=0$, and
$\mathrm c_0I_K+\mathbf M^\star=\exp\bigl((\log\mathrm c_0)I_K+\mathrm c_0^{-1}\mathbf M^\star\bigr)$. Then
\eqref{eq:link-strong-convexity} gives the infimum and the convergence of minimizing sequences, and the
infimum is not attained because chart links are invertible and $\mathbf M^\star$ is not. If
$\mu_j^\top\Sigma_j^{-1}\mu_i<0$, then $\operatorname{tr}\mathbf M^\star<0$ is a simple eigenvalue of $\mathbf M^\star$; for
$\mathbf M$ near $\mathbf M^\star$ exactly one eigenvalue lies in the disk of radius
$\frac12\lvert\operatorname{tr}\mathbf M^\star\rvert$ about $\operatorname{tr}\mathbf M^\star$, by continuity of the roots of
the characteristic polynomial, and it is real, negative and simple because the disk is symmetric under
conjugation. A real matrix $e^{\mathbf B}$ has no simple negative eigenvalue: triangularizing $\mathbf B$ over
$\mathbb C$ shows that the eigenvalues of $e^{\mathbf B}$ with multiplicity are the $e^\beta$ for the eigenvalues
$\beta$ of $\mathbf B$, a negative value requires $\beta\notin\R$, and then $\bar\beta\neq\beta$ is an eigenvalue with
the same image; this is a case of the criterion of \citet[Theorem~1]{culver1966existence}: a real
invertible matrix has a real logarithm exactly when every Jordan block belonging to a negative eigenvalue
occurs an even number of times. So a neighborhood of $\mathbf M^\star$ contains no chart link, and
\eqref{eq:link-strong-convexity} bounds $\mathbf l$ on its complement away from $\mathbf l(\mathbf M^\star)$. For
$K=2$, the chart links are the matrices of $\exp(\mathfrak{gl}(2,\R))$. By \eqref{eq:link-strong-convexity},
$\mathbf l$ attains its minimum over the closure of this set, that minimum is the infimum over chart links,
and a chart link attaining the infimum would minimize $\mathbf l$ over the closure. No chart link $\mathbf M$
does, because such an $\mathbf M$ would satisfy $\nabla\mathbf l(\mathbf M)=0$, that is $\mathbf M=\mathbf M^\star$,
which is not in the closure. If $\mathbf M$ has no eigenvalue in $(-\infty,0]$, then neither has any matrix
near it, and these matrices are chart links by the criterion, so $\mathbf M$ is an interior point of the
closure and the gradient of $\mathbf l$ vanishes there. Otherwise the criterion leaves only $\mathbf M=\mathrm c'I_2$ with $\mathrm c'<0$, since the second
eigenvalue of a real $2\times2$ matrix with a negative eigenvalue is real. For every $\mathbf E$ with
$(\mathbf E_{11}-\mathbf E_{22})^2+4\mathbf E_{12}\mathbf E_{21}<0$ and every $t>0$ the matrix $\mathrm c'I_2+t\mathbf E$ has
non-real eigenvalues and is a chart link, so minimality gives
$\langle\nabla\mathbf l(\mathrm c'I_2),\mathbf E\rangle\geq0$; these $\mathbf E$ form a nonempty open set invariant
under $\mathbf E\mapsto-\mathbf E$, so $\nabla\mathbf l(\mathrm c'I_2)$ is orthogonal to an open set and vanishes.
(b) This is \cref{prop:link-law-level-energy}(a,b), since
$(K-1)(K-2)/2\geq1$ for $K\geq3$. (c)(i) $\mathbf V$ is continuous and at least
$\frac12\langle\mathring\omega-\mathring\alpha,\mathring\Sigma^{-1}(\mathring\omega-\mathring\alpha)\rangle+\inf\mathbf U$, which tends to
$+\infty$ with $\lVert\mathring\omega\rVert$. By \eqref{eq:link-field-free-energy} the free energy is at least
$\KL(R\Vert\nu^{\mathscr T})+\inf\mathbf U+\log Z_{\mathscr S}+\E_{Q^{\mathsf Y}}\log q^{\mathsf Y}$. With $\mathring\lambda_l$ the
eigenvalues of $\mathring\Sigma^{-1/2}\mathring\Lambda^{-1}\mathring\Sigma^{-1/2}$, \eqref{eq:link-gaussian-kl} is
$\frac12\sum_l(\mathring\lambda_l-1-\log\mathring\lambda_l)+\frac12\langle\mathring\mu-\mathring\alpha,\mathring\Sigma^{-1}(\mathring\mu-\mathring\alpha)\rangle$,
so a bound on it confines the $\mathring\lambda_l$ to a compact subset of $(0,\infty)$ and $\mathring\mu$ to a bounded set:
the sublevel sets are compact in the open domain, and the free energy is continuous there by
\cref{prop:link-field-equation}(a). (ii) The record densities are at most one and positive, and by
\cref{prop:link-gaussian-recognition-vfe}(c) the free energy is
$\KL(R\Vert\mathbb P^{\mathscr T}_\theta)+\KL(Q^{\mathsf Y}\Vert\mathbb P^{\mathsf Y})+\E_{R\otimes Q^{\mathsf Y}}H_o$, finite everywhere by the
growth bound, and at least $\KL(R\Vert\nu^{\mathscr T})+\log Z_{\mathscr S}+\KL(Q^{\mathsf Y}\Vert\mathbb P^{\mathsf Y})$ because
$\mathscr S\geq0$ and $H_o\geq0$; both divergences have compact sublevel sets in their parameters by the
argument of (i). The divergences and $\E_R\mathscr S$ are continuous, and $\E_{R\otimes Q^{\mathsf Y}}H_o$ is lower
semicontinuous by Fatou's lemma in the representation by standard Gaussian vectors, so the minimum is
attained. \hfill$\square$

\noindent\emph{Proof of \Cref{thm:link-held-limit}.}
Write $\mathring\omega=\mathring\alpha+\mathring\sigma\mathring\Sigma_0^{1/2}\mathbf b$
under $\nu^{\mathscr T}_{\mathring\sigma}$, with $\mathbf b$ standard Gaussian. Since $0\leq e^{-\mathscr S}\leq1$ is
continuous, $Z_{\mathscr S,\mathring\sigma}=\E e^{-\mathscr S(\mathring\alpha+\mathring\sigma\mathring\Sigma_0^{1/2}\mathbf b)}\to e^{-\mathscr S(\mathring\alpha)}>0$,
and
$\mathbb P^{\mathscr T}_{\mathring\sigma}\{\lVert\mathring\omega-\mathring\alpha\rVert\geq\mathring\epsilon\}\leq\nu^{\mathscr T}_{\mathring\sigma}\{\lVert\mathring\omega-\mathring\alpha\rVert\geq\mathring\epsilon\}/Z_{\mathscr S,\mathring\sigma}\to0$
for every $\mathring\epsilon>0$. By \eqref{eq:link-decomposition}, every law $Q=R(d\mathring\omega)Q_{\mathring\omega}(dy)$ has
$\Fenergy_{\mathring\sigma}[Q]=\KL(R\Vert\mathbb P^{\mathscr T}_{\mathring\sigma})+\int\mathbf F_{Q_{\mathring\omega}}(\mathring\omega)R(d\mathring\omega)$, with
$\mathbf F_{Q_{\mathring\omega}}\geq-\log p(o\given\cdot)\geq-\log\mathrm c_{\mathfrak A}$.

Upper bound. Let $Q\in\mathcal Q'_0$ with $\mathbf F_Q(\mathring\alpha)<\infty$. The law
$\nu^{\mathscr T}_{\mathring\sigma}\otimes Q$ is admissible, and
$\KL(\nu^{\mathscr T}_{\mathring\sigma}\Vert\mathbb P^{\mathscr T}_{\mathring\sigma})=\E_{\nu^{\mathscr T}_{\mathring\sigma}}\mathscr S+\log Z_{\mathscr S,\mathring\sigma}$.
For $\mathring\sigma\leq1$, \cref{lem:link-growth}(c) dominates
$\mathscr S(\mathring\alpha+\mathring\sigma\mathring\Sigma_0^{1/2}\mathbf b)$ by
$\mathrm c_1e^{\mathrm c_2(\lVert\mathring\alpha\rVert+\lVert\mathring\Sigma_0\rVert^{1/2}\lVert\mathbf b\rVert)}$, which is integrable by
\cref{lem:link-growth}(a), so $\E_{\nu^{\mathscr T}_{\mathring\sigma}}\mathscr S\to\mathscr S(\mathring\alpha)$ and
$\KL(\nu^{\mathscr T}_{\mathring\sigma}\Vert\mathbb P^{\mathscr T}_{\mathring\sigma})\to0$. Hypothesis (iii) and the same domination
give $\E_{\nu^{\mathscr T}_{\mathring\sigma}}\mathbf F_Q\to\mathbf F_Q(\mathring\alpha)$. Hence the product infimum has
$\limsup$ at most $\mathbf F_Q(\mathring\alpha)$, and taking the infimum over $\mathcal Q'_0$ gives
$\limsup\leq\mathbf F(\mathring\alpha)$.

Lower bound. Let $\mathrm c_0>0$. By (ii) there is $\mathring\epsilon>0$ with $\mathbf F>\mathbf F(\mathring\alpha)-\mathrm c_0$
on $\{\lVert\mathring\omega-\mathring\alpha\rVert<\mathring\epsilon\}$, so the bounded simple function
$\mathbf q=(\mathbf F(\mathring\alpha)-\mathrm c_0)\mathbf 1_{\{\lVert\mathring\omega-\mathring\alpha\rVert<\mathring\epsilon\}}-(\log\mathrm c_{\mathfrak A})\mathbf 1_{\{\lVert\mathring\omega-\mathring\alpha\rVert\geq\mathring\epsilon\}}$
satisfies
$\mathbf q\leq\mathbf F_{Q_{\mathring\omega}}$ for every kernel with values in $\mathcal Q'$. For every law $R$ with
$\KL(R\Vert\mathbb P^{\mathscr T}_{\mathring\sigma})<\infty$, the law
$\mathbb P^{\mathbf q}\propto e^{-\mathbf q}\mathbb P^{\mathscr T}_{\mathring\sigma}$ has a density bounded above and below, so
$0\leq\KL(R\Vert\mathbb P^{\mathbf q})=\KL(R\Vert\mathbb P^{\mathscr T}_{\mathring\sigma})+\E_R\mathbf q+\log\E_{\mathbb P^{\mathscr T}_{\mathring\sigma}}e^{-\mathbf q}$,
and therefore
$\Fenergy_{\mathring\sigma}[Q]\geq\KL(R\Vert\mathbb P^{\mathscr T}_{\mathring\sigma})+\E_R\mathbf q\geq-\log\E_{\mathbb P^{\mathscr T}_{\mathring\sigma}}e^{-\mathbf q}$;
the bound holds trivially when $\KL(R\Vert\mathbb P^{\mathscr T}_{\mathring\sigma})=+\infty$. Since
$\mathbb P^{\mathscr T}_{\mathring\sigma}\{\lVert\mathring\omega-\mathring\alpha\rVert\geq\mathring\epsilon\}\to0$, the right side tends to
$\mathbf F(\mathring\alpha)-\mathrm c_0$. The
kernel infimum is at most the product infimum, so both converge to $\mathbf F(\mathring\alpha)$.

(b) The decomposition gives
$\KL(R_{\mathring\sigma}\Vert\mathbb P^{\mathscr T}_{\mathring\sigma})\leq\Fenergy_{\mathring\sigma}+\log\mathrm c_{\mathfrak A}\leq\mathrm c$. Data processing
(\cref{lem:prob-kl-chain-rule}(d)) through the indicator of $\{\lVert\mathring\omega-\mathring\alpha\rVert\geq\mathring\epsilon\}$ bounds
$\mathrm c$ below by the relative entropy between the Bernoulli laws with success probabilities
$R_{\mathring\sigma}\{\lVert\mathring\omega-\mathring\alpha\rVert\geq\mathring\epsilon\}$ and
$\mathbb P^{\mathscr T}_{\mathring\sigma}\{\lVert\mathring\omega-\mathring\alpha\rVert\geq\mathring\epsilon\}$. Since the entropy of a Bernoulli law is
at most $\log2$ and the failure outcome contributes a nonnegative cross-entropy term, that relative
entropy is at least
$R_{\mathring\sigma}\{\lVert\mathring\omega-\mathring\alpha\rVert\geq\mathring\epsilon\}\log\bigl(1/\mathbb P^{\mathscr T}_{\mathring\sigma}\{\lVert\mathring\omega-\mathring\alpha\rVert\geq\mathring\epsilon\}\bigr)-\log2$,
so $R_{\mathring\sigma}\{\lVert\mathring\omega-\mathring\alpha\rVert\geq\mathring\epsilon\}\leq(\mathrm c+\log2)/\log\bigl(1/\mathbb P^{\mathscr T}_{\mathring\sigma}\{\lVert\mathring\omega-\mathring\alpha\rVert\geq\mathring\epsilon\}\bigr)\to0$.
For the second claim, with
$\mathbf q$ as above,
$\KL(R_{\mathring\sigma}\Vert\mathbb P^{\mathscr T}_{\mathring\sigma})=\Fenergy_{\mathring\sigma}-\int\mathbf F_{Q_{\mathring\sigma,\mathring\omega}}dR_{\mathring\sigma}\leq\Fenergy_{\mathring\sigma}-\E_{R_{\mathring\sigma}}\mathbf q$,
and $\E_{R_{\mathring\sigma}}\mathbf q\to\mathbf F(\mathring\alpha)-\mathrm c_0$ by the first claim, so the $\limsup$ of the
divergence is at most $\mathrm c_0$.

(c) By \eqref{eq:link-latent-evidence},
$\lvert p_{\theta,V}(o\given X)-p(o\given\mathring\alpha)\rvert\leq\sup_{\lVert\mathring\omega-\mathring\alpha\rVert<\mathring\epsilon}\lvert p(o\given\cdot)-p(o\given\mathring\alpha)\rvert
+\mathrm c_{\mathfrak A}\mathbb P^{\mathscr T}_{\mathring\sigma}\{\lVert\mathring\omega-\mathring\alpha\rVert\geq\mathring\epsilon\}$,
which tends to zero when $\mathring\epsilon$ is chosen first. Then
$\boldsymbol\Pi^{\mathscr T}_{\mathring\sigma}\{\lVert\mathring\omega-\mathring\alpha\rVert\geq\mathring\epsilon\}\leq\mathrm c_{\mathfrak A}\mathbb P^{\mathscr T}_{\mathring\sigma}\{\lVert\mathring\omega-\mathring\alpha\rVert\geq\mathring\epsilon\}/p_{\theta,V}(o\given X)\to0$,
because the evidence tends to $p(o\given\mathring\alpha)>0$. For the class of all laws on $\mathsf Y$ the
infimum of the free energy is minus the log evidence (\cref{thm:elbo-extended-gap}), so the first claim is
the value statement of (a) for that class, here obtained without (ii) and (iii).

(d) Suppose $\mathring\omega_l\to\mathring\alpha$ with $\mathbf F(\mathring\omega_l)\to\mathrm c'<\mathbf F(\mathring\alpha)$ along a
subsequence. For large $l$ choose $\lambda_l$ with
$\mathbf F_{Q_{\lambda_l}}(\mathring\omega_l)\leq\min\{\mathbf F(\mathring\omega_l)+1/l,\mathrm c\}$; the $\lambda_l$ lie in a compact
subset of $\mathbb U$, a further subsequence converges to some $\lambda\in\mathbb U$, and lower semicontinuity
gives $\mathbf F_{Q_\lambda}(\mathring\alpha)\leq\mathrm c'<\mathbf F(\mathring\alpha)$, which contradicts the definition of $\mathbf F$.
\hfill$\square$

\noindent\emph{Proof of \Cref{prop:link-recognition-flow}.}
For \eqref{eq:link-gaussian-free-energy}, the link sector is
$\KL(R\Vert\nu^{\mathscr T})+\E_R\mathscr S+\log Z_{\mathscr S}$ by \eqref{eq:link-sector}. With $r$ the
Lebesgue density of $R$,
$\KL(R\Vert\nu^{\mathscr T})=\E_R[\log r]-\E_R[\log(d\nu^{\mathscr T}/d\mathring\omega)]$, where
$\E_R[\log r]=\tfrac12\log\det\mathring\Lambda-\tfrac12\dim\mathscr T\log(2\pi e)$ and
$-\log(d\nu^{\mathscr T}/d\mathring\omega)$ is the sum of the last two terms of
\eqref{eq:link-generator-integrand}; substituting in \eqref{eq:link-decomposition} gives the display.

(a) The density of $R$ is
$\exp(\langle\mathring\vartheta_1,\mathring\omega\rangle+\langle\mathring\omega,\mathring\vartheta_2\mathring\omega\rangle-\mathring\Gamma(\mathring\vartheta))$
with
$\mathring\Gamma(\mathring\vartheta)=-\tfrac14\langle\mathring\vartheta_1,\mathring\vartheta_2^{-1}\mathring\vartheta_1\rangle-\tfrac12\log\det(-2\mathring\vartheta_2)+\tfrac12\dim\mathscr T\log(2\pi)$.
For $\mathring\vartheta_2\in\mathring\Xi$,
$\langle\mathring\omega,\mathring\vartheta_2\mathring\omega\rangle=\Tr(\mathring\vartheta_2\operatorname{pr}_{\mathring\Xi}(\mathring\omega\otimes\mathring\omega))$,
so the sufficient statistic is
$(\mathring\omega,\operatorname{pr}_{\mathring\Xi}(\mathring\omega\otimes\mathring\omega))$.
The integral of $\exp(\langle\mathring\vartheta_1,\mathring\omega\rangle+\langle\mathring\omega,\mathring\vartheta_2\mathring\omega\rangle)$
is finite exactly when $\mathring\vartheta_2$ is negative definite, an open condition in
$\mathscr T\times\mathring\Xi$, so the family is regular, and it is minimal because a nonzero polynomial
of degree at most two in $\mathring\omega$ vanishes only on a Lebesgue-null set. A positive-definite
operator that preserves each summand has an inverse that preserves each summand, so
$\mathring\Lambda^{-1}\in\mathring\Xi$, and $\nabla\mathring\Gamma$, the $R$-mean of the sufficient
statistic, is $\mathring\varpi$ as stated. Under $R$ the centered generator
$\mathring\omega-\mathring\mu$ has law $\mathcal N(0,\mathring\Lambda^{-1})$, and the score of $R$ in the
direction $(\delta\mathring\mu,\delta\mathring\Lambda)$ is
$\langle\delta\mathring\mu,\mathring\Lambda(\mathring\omega-\mathring\mu)\rangle-\tfrac12(\langle\mathring\omega-\mathring\mu,\delta\mathring\Lambda(\mathring\omega-\mathring\mu)\rangle-\Tr(\mathring\Lambda^{-1}\delta\mathring\Lambda))$.
The two parts are uncorrelated because the odd central moments vanish, the first has variance
$\langle\delta\mathring\mu,\mathring\Lambda\delta\mathring\mu\rangle$, and Isserlis' identity, by which
$\langle\mathring\omega-\mathring\mu,\mathring\Delta(\mathring\omega-\mathring\mu)\rangle$ has variance
$2\Tr(\mathring\Delta\mathring\Lambda^{-1}\mathring\Delta\mathring\Lambda^{-1})$ for self-adjoint $\mathring\Delta$, gives the second part of
\eqref{eq:link-gaussian-fisher}. The form vanishes only
when $\delta\mathring\mu=0$ and $\mathring\Lambda^{-1/2}\delta\mathring\Lambda\mathring\Lambda^{-1/2}=0$,
and its coefficients are smooth. Every minimal regular exponential family with its Fisher metric is an
instance of \cref{hyp:hist-standing-configuration-tier}, and a finite product of instances with the
direct-sum metric, $C^1$ and positive definite, is again one.

(b) The metric $\mathsf G$ is $C^1$ and $F$ is $C^2$, so $\operatorname{grad}^{\mathsf G}F$ is $C^1$
and the Picard--Lindel\"of theorem gives existence and uniqueness of maximal integral curves. Along
one, $\tfrac{d}{dt}F=dF[\dot Q]=-\mathsf G(\operatorname{grad}^{\mathsf G}F,\operatorname{grad}^{\mathsf G}F)$
by \eqref{eq:hist-fisher-gradient-definition}. By \cref{prop:hist-vfe-orbit-embedded}, applied with
$\mathsf M=\mathsf G^{-1}$, a curve through a noncritical point meets no critical point, so its
velocity never vanishes, and it satisfies \eqref{eq:hist-vfe-ray-representative} with
$\mathsf a\equiv1$.

(c) The field depends on the configuration only and its zeros are the zeros of $dF$. The force
balance is \eqref{eq:hist-fisher-gradient-definition} with $Z$ arbitrary and
$\dot Q=-\operatorname{grad}^{\mathsf G}F$, and evaluating it on $\dot Q$ gives the dissipation.
Strict decrease along nonconstant solutions is \cref{prop:hist-vfe-orbit-embedded}, and a periodic
nonconstant solution would return to a value of $F$ it has left. In a chart, the derivative of
$Q\mapsto-\mathsf G(Q)^{-1}\nabla F(Q)$ at $Q_\star$ is $-\mathsf G(Q_\star)^{-1}\nabla^2F(Q_\star)$,
because the term containing the derivative of $\mathsf G^{-1}$ is multiplied by
$\nabla F(Q_\star)=0$; this matrix is similar to the symmetric matrix
$-\mathsf G(Q_\star)^{-1/2}\nabla^2F(Q_\star)\mathsf G(Q_\star)^{-1/2}$, hence diagonalizable with real
eigenvalues, and every solution of the linearized equation is a combination of real exponentials.

(d) In the natural chart the Fisher metric is $\nabla^2\mathring\Gamma$, and the chain rule through
$\mathring\varpi=\nabla\mathring\Gamma(\mathring\vartheta)$ gives
$\nabla_{\mathring\vartheta}F=\nabla^2\mathring\Gamma\nabla_{\mathring\varpi}F$, so the generator block of
the flow is $\dot{\mathring\vartheta}=-\nabla_{\mathring\varpi}F$ and its Euler step is
$\mathring\vartheta'=\mathring\vartheta-\Delta t\nabla_{\mathring\varpi}F$. The growth bound and the
Gaussian tails of $R$ justify differentiation under the integral and two integrations by parts, which
give the Gaussian gradient identities of \cref{lem:link-gaussian-gradients}:
$\nabla_{\mathring\mu}\E_R[\mathfrak i_\lambda]=\E_R[\nabla\mathfrak i_\lambda]$, and the derivative of
$\E_R[\mathfrak i_\lambda]$ in the covariance $\mathring\Lambda^{-1}$ along $\mathring\Delta\in\mathring\Xi$ is
$\tfrac12\Tr(\mathring\Delta\E_R[\nabla^2\mathfrak i_\lambda])$, so its gradient in $\mathring\Xi$ is
$\tfrac12\operatorname{pr}_{\mathring\Xi}\E_R[\nabla^2\mathfrak i_\lambda]$. The covariance gradient of
$\tfrac12\log\det\mathring\Lambda=-\tfrac12\log\det\mathring\Lambda^{-1}$ is
$-\tfrac12\mathring\Lambda$. Write $\mathring\Upsilon=\tfrac12\operatorname{pr}_{\mathring\Xi}\E_R[\nabla^2\mathfrak i_\lambda]-\tfrac12\mathring\Lambda\in\mathring\Xi$
for the covariance gradient of $F$. Since $\mathring\mu=\mathring\varpi_1$ and
$\mathring\Lambda^{-1}=\mathring\varpi_2-\operatorname{pr}_{\mathring\Xi}(\mathring\varpi_1\otimes\mathring\varpi_1)$,
and $\Tr(\mathring\Upsilon\operatorname{pr}_{\mathring\Xi}(\mathring\Delta))=\Tr(\mathring\Upsilon\mathring\Delta)$
for every self-adjoint $\mathring\Delta$, the chain rule gives $\nabla_{\mathring\varpi_2}F=\mathring\Upsilon$ and
$\nabla_{\mathring\varpi_1}F=\E_R[\nabla\mathfrak i_\lambda]-2\mathring\Upsilon\mathring\mu$. The second component
of the Euler step is
$-\tfrac12\mathring\Lambda'=-\tfrac12\mathring\Lambda-\Delta t\mathring\Upsilon$, which is the first equation
of \eqref{eq:link-blr}, and the first component is
$\mathring\Lambda'\mathring\mu'=\mathring\Lambda\mathring\mu-\Delta t\E_R[\nabla\mathfrak i_\lambda]+2\Delta t\mathring\Upsilon\mathring\mu=\mathring\Lambda'\mathring\mu-\Delta t\E_R[\nabla\mathfrak i_\lambda]$,
which is the second. When $0<\Delta t<1$ and the projected expected Hessian is positive semidefinite,
$\mathring\Lambda'\succeq(1-\Delta t)\mathring\Lambda$ is positive definite. For one summand
$\operatorname{pr}_{\mathring\Xi}$ is the identity and \eqref{eq:link-blr} has the form of
\citet[Eq.~(12)]{KhanRue2023} with their loss replaced by $\mathfrak i_\lambda$, and for the lines of an
orthonormal basis $\operatorname{pr}_{\mathring\Xi}$ keeps the diagonal, as in their Eq.~(37), where the
expected gradient and Hessian are evaluated at the mean.

(e) The Lebesgue density of $R$ gives
$\E_R[\log r]=\langle\mathring\vartheta,\mathring\varpi\rangle-\mathring\Gamma(\mathring\vartheta)=\mathring\Gamma^\star(\mathring\varpi)$,
so $\widetilde F=\E_R[\mathfrak i_\lambda]+\mathring\Gamma^\star+\log Z_{\mathscr S}$. The relative entropy of
two members is
$\KL(R\Vert R')=\mathring\Gamma(\mathring\vartheta')-\mathring\Gamma(\mathring\vartheta)-\langle\mathring\varpi,\mathring\vartheta'-\mathring\vartheta\rangle$,
and the Bregman divergence of $\mathring\Gamma^\star$ is the relative entropy with the arguments
exchanged,
$\mathring\Gamma^\star(\mathring\varpi')-\mathring\Gamma^\star(\mathring\varpi)-\langle\mathring\vartheta,\mathring\varpi'-\mathring\varpi\rangle=\KL(R'\Vert R)$.
The expectation domain $\{(\mathring\varpi_1,\mathring\varpi_2):\mathring\varpi_2-\operatorname{pr}_{\mathring\Xi}(\mathring\varpi_1\otimes\mathring\varpi_1)\succ0\}$
is convex, because along a segment the difference of convex combinations is the convex combination of
the differences plus a nonnegative multiple of $\operatorname{pr}_{\mathring\Xi}$ of a positive-semidefinite
operator, which is positive semidefinite because $\operatorname{pr}_{\mathring\Xi}$ keeps the diagonal
blocks of the decomposition. Taylor's formula with
integral remainder along the segment and the curvature hypothesis give
$\widetilde F(\mathring\varpi')\leq\widetilde F(\mathring\varpi)+\langle\nabla\widetilde F(\mathring\varpi),\mathring\varpi'-\mathring\varpi\rangle+\mathrm c_{\mathsf G}\KL(R'\Vert R)$,
since the same formula applied to $\mathring\Gamma^\star$ produces the Bregman divergence. By (d),
$\nabla\widetilde F(\mathring\varpi)=\Delta t^{-1}(\mathring\vartheta-\mathring\vartheta')$, and adding the two
expressions for the divergences gives
$\langle\mathring\vartheta'-\mathring\vartheta,\mathring\varpi'-\mathring\varpi\rangle=\KL(R\Vert R')+\KL(R'\Vert R)$;
substituting proves \eqref{eq:link-blr-descent}. Both divergences vanish only when $R'=R$. For the
quadratic integrand,
$\E_R[\mathfrak i_\lambda]=\tfrac12\Tr(\mathbf Q\mathring\varpi_2)-\langle\mathbf b,\mathring\varpi_1\rangle+\mathrm c_3$
is affine in $\mathring\varpi$, so $\nabla^2\widetilde F=\nabla^2\mathring\Gamma^\star$; moreover
$\E_R[\nabla^2\mathfrak i_\lambda]=\mathbf Q$ and $\E_R[\nabla\mathfrak i_\lambda]=\mathbf Q\mathring\mu-\mathbf b$, so
\eqref{eq:link-blr} with $\Delta t=1$ gives $\mathring\Lambda'=\mathbf Q$ and
$\mathring\mu'=\mathbf Q^{-1}\mathbf b$, and $R'$ is the law proportional to $e^{-\mathfrak i_\lambda}$, which
minimizes $F=\KL(R\Vert R')+\mathrm{const}$ over $R$. \hfill$\square$

\noindent\emph{Proof of \Cref{prop:link-anchor-m-coordinate}.}
(a) The joint log density is the prior log density plus terms free of $\mathring\alpha$, and
$\log(d\mathbb P^{\mathscr T}_{\mathring\alpha}/d\mathring\omega)=-\mathscr S(\mathring\omega)-\log Z_{\mathscr S}(\mathring\alpha)-\tfrac12\langle\mathring\omega-\mathring\alpha,\mathring\Sigma^{-1}(\mathring\omega-\mathring\alpha)\rangle-\tfrac12\log\det(2\pi\mathring\Sigma)$,
which gives \eqref{eq:link-anchor-objective}. Differentiating under the integral, which the factor
$e^{-\mathscr S}\leq1$ and the Gaussian tails permit locally uniformly in $\mathring\alpha$,
$\nabla_{\mathring\alpha_e}\log Z_{\mathscr S}=\mathring\Sigma_e^{-1}(\E_{\mathbb P^{\mathscr T}_{\mathring\alpha}}[\mathring\omega_e]-\mathring\alpha_e)$,
and subtracting it from the gradient of the first term of \eqref{eq:link-anchor-objective} gives
\eqref{eq:link-anchor-gradient}.

(b) With $\mathring\pi=\mathring\Sigma^{-1}\mathring\alpha$, the prior density with respect to the
measure
$\exp(-\tfrac12\langle\mathring\omega,\mathring\Sigma^{-1}\mathring\omega\rangle-\mathscr S(\mathring\omega))d\mathring\omega$
is $\exp(\langle\mathring\pi,\mathring\omega\rangle-\mathrm A(\mathring\pi))$, with
$\mathrm A(\mathring\pi)=\log\int\exp(\langle\mathring\pi,\mathring\omega\rangle-\tfrac12\langle\mathring\omega,\mathring\Sigma^{-1}\mathring\omega\rangle-\mathscr S(\mathring\omega))d\mathring\omega$,
finite on all of $\mathscr T$ because $\mathscr S\geq0$. Thus
$\Lelbo^{\mathrm{anc}}=\langle\mathring\pi,\E_R[\mathring\omega]\rangle-\mathrm A(\mathring\pi)+\mathrm{const}$, and
$\mathrm A$ is smooth with Hessian $\operatorname{Cov}_{\mathbb P^{\mathscr T}_{\mathring\alpha}}(\mathring\omega)$,
which is positive definite because the prior has a positive density; so $\Lelbo^{\mathrm{anc}}$ is strictly concave in
$\mathring\pi$, hence in $\mathring\alpha$. For a unit vector $\mathring\delta$, the set
$\{\langle\mathring\delta,\mathring\omega\rangle\geq\langle\mathring\delta,\E_R[\mathring\omega]\rangle+1\}$ has
positive measure $\mathrm c_{\mathring\delta}$ under that measure, so the convex function
$\mathrm f=\mathrm A-\langle\cdot,\E_R[\mathring\omega]\rangle$ satisfies
$\mathrm f(t\mathring\delta)\geq t+\log\mathrm c_{\mathring\delta}\to\infty$ as $t\to\infty$. A convex function
on $\mathscr T$ that tends to infinity along every ray from the origin has bounded sublevel sets: if
$\mathrm f(\mathring\pi_{\mathrm l})\leq\mathrm c$ with $\lVert\mathring\pi_{\mathrm l}\rVert\to\infty$ and, after
passing to a subsequence, $\mathring\pi_{\mathrm l}/\lVert\mathring\pi_{\mathrm l}\rVert\to\mathring\delta$,
convexity gives
$\mathrm f(t\mathring\pi_{\mathrm l}/\lVert\mathring\pi_{\mathrm l}\rVert)\leq\max(\mathrm f(0),\mathrm c)$ for
$\lVert\mathring\pi_{\mathrm l}\rVert\geq t$, and continuity bounds $\mathrm f(t\mathring\delta)$ for every $t$. Hence
$\mathrm f$ attains its minimum, at exactly one point, where
$\nabla\mathrm A=\E_{\mathbb P^{\mathscr T}_{\mathring\alpha}}[\mathring\omega]=\E_R[\mathring\omega]$.
Since the other factors of the joint do not depend on $\mathring\alpha$, the maximizers of
\eqref{eq:elbo-m-coordinate} over the anchor are those of $\Lelbo^{\mathrm{anc}}$, and \cref{prop:elbo-exact-m-coordinate}
applies.

(c) If $\mathscr S=0$ the prior is $\nu^{\mathscr T}_{\mathring\alpha}$, with mean $\mathring\alpha$,
and the moment condition reads $\mathring\alpha=\E_R[\mathring\omega]$. In the affine representation
$\rho(\mathring\omega)=\left(\begin{smallmatrix}1&\mathring\omega\\0&1\end{smallmatrix}\right)$ of the
translations, products of represented translations are represented sums, so the represented
holonomy of $p$ minus the identity is $\mathrm d\rho(\mathbf C_p\mathring\omega)$ and the loop defect
is exactly $\tfrac12\lVert\mathbf C_p\mathring\omega\rVert_{\mathsf G}^2$, as in
\cref{prop:link-translation-instance}. For a quadratic action, the
prior is Gaussian with precision $\mathring\Sigma^{-1}+\sum_p\mathrm c_p\mathbf C_p^\top\mathbf C_p$ and
mean $(\mathring\Sigma^{-1}+\sum_p\mathrm c_p\mathbf C_p^\top\mathbf C_p)^{-1}\mathring\Sigma^{-1}\mathring\alpha$,
and the moment condition gives \eqref{eq:link-anchor-quadratic}. The operator
$\sum_p\mathrm c_p\mathbf C_p^\top\mathbf C_p$ is positive semidefinite with quadratic form
$\sum_p\mathrm c_p\lVert\mathbf C_p\mathring\omega\rVert_{\mathsf G}^2$, so it annihilates
$\E_R[\mathring\omega]$ exactly when every $\mathbf C_p\E_R[\mathring\omega]$ with $\mathrm c_p>0$
vanishes, and $\mathring\Sigma$ is invertible.

(d) In the coordinates $(\mathring\Sigma_e^{-1}\mathring\alpha_e,-\tfrac12\mathring\Sigma_e^{-1})_{e\in E_+}$
the prior is the exponential family with sufficient statistic
$(\mathring\omega_e,\mathring\omega_e\otimes\mathring\omega_e)_{e\in E_+}$ and base measure
$e^{-\mathscr S(\mathring\omega)}d\mathring\omega$; the Gaussian normalizers of the base law and
$Z_{\mathscr S}$ combine into its log-partition function, which is strictly convex because the
components of the statistic are affinely independent on a set of positive Lebesgue measure. The
range of these coordinates is convex, so the strictly concave objective has at most one maximizer,
and at an interior maximizer the gradient of the log-partition function, the vector of prior moments
of the statistic, equals its expectation under $R$, which is the stated matching. If $\mathscr S=0$,
the prior is $\bigotimes_e\mathcal N(\mathring\alpha_e,\mathring\Sigma_e)$, the matching is solved by
the stated pair, which lies in the range when every $\operatorname{Cov}_R(\mathring\omega_e)$ is
positive definite, and a critical point of a concave function maximizes it. For the witness, with
$\mathsf G^{\mathrm{fib}}=I_2$ the inner product on $\mathfrak g=\R$ is the product of reals, the loop
defect is $\tfrac12\mathring\omega_e^2$, and with $\mathrm s=\mathring\Sigma_e$ the prior is
$\mathcal N(\mathring\alpha_e/(1+\mathrm c\mathrm s),\mathrm s/(1+\mathrm c\mathrm s))$. Its variance
$\mathrm s'=\mathrm s/(1+\mathrm c\mathrm s)$ ranges over $(0,1/\mathrm c)$, and for each $\mathrm s$
the best anchor makes the prior mean $\mathring\mu$, leaving
$-\tfrac12\log(2\pi\mathrm s')-\mathrm v/(2\mathrm s')$, whose derivative in $\mathrm s'$ is
$(\mathrm v-\mathrm s')/(2\mathrm s'^2)>0$ for $\mathrm s'<1/\mathrm c\leq\mathrm v$. The objective
therefore increases strictly as $\mathrm s\to\infty$ and no pair attains its supremum.
\hfill$\square$

\noindent\emph{Proof of \Cref{rem:link-rates}.}
(i) A frame at the reference point is $U_i=e_G$, and then $\mathring\Theta_e=\Theta_e$ by
\eqref{eq:geo-edge-twist}. In the affine representation a product of represented translations is the
represented sum, so a represented holonomy minus the identity is $\mathrm d\rho$ of the loop sum, and
for $\mathsf G^{\mathrm{fib}}=I_2$,
$\langle\mathrm d\rho\mathrm x,\mathrm d\rho\mathrm y\rangle_{\mathsf G}=\Tr\bigl((\mathrm xE_{12})^\top\mathrm yE_{12}\bigr)=\mathrm x\mathrm y$.
The configuration prior of \cref{def:rg-formation-cycle-experiment} is $\mathcal N(0,I_{12})$, whose
$\check\omega$-block is $\mathcal N(0,I_4)$ and independent of $(\bar k,\bar m)$, and the records are
the binary records with success probability $\tfrac12\exp[-\mathcal D_{e(i,S_i)}/\mathrm T]$, whose
edge costs carry the loop structure. Disintegrating a configuration law over $\check\omega$ gives the
form $R(d\check\omega)Q_{\check\omega}$ of \cref{prop:elbo-random-link-vfe}. (ii) is
\cref{prop:obs-slow-model-link-dichotomy}(i) with the link block of the mean read as
$\mathring\mu$ by (i). (iii) is \cref{def:obs-fast-slow-blocks}, whose slow coordinate carries
$\theta$, together with the statement of \cref{def:rg-formation-cycle-experiment} that no generative
parameter is learned.

(iv) Write $U=U_{S^{(0)}}$. It is the quadratic form
$U(\mathbf x)=\tfrac12\mathbf x^\top\nabla^2U\mathbf x$, and $\nabla^2U\succeq\mathrm TI$ because
$U-\tfrac{\mathrm T}2\lvert\mathbf x\rvert^2$ is a sum of squares. By
\cref{prop:rg-formation-endogenous-cycle}, $\Fenergy_{\mathrm T}(q)=\mathrm T\KL(q\Vert q^\star)+\mathrm{const}$
with $q^\star\propto\pi_Se^{-U_S/\mathrm T}$, so at $q=\mathcal N(\mu,\Sigma)\otimes\delta_{S^{(0)}}$,
$F=\Fenergy_{\mathrm T}/\mathrm T=\mathrm T^{-1}U(\mu)+\tfrac12\mathrm T^{-1}\Tr(\nabla^2U\Sigma)-\tfrac12\log\det\Sigma+\mathrm{const}$.
In covariance coordinates, where $\delta\Lambda=-\Lambda\delta\Sigma\Lambda$ turns
\eqref{eq:link-gaussian-fisher} into
$\langle\delta\mu,\Sigma^{-1}\delta\mu\rangle+\tfrac12\Tr(\Sigma^{-1}\delta\Sigma\Sigma^{-1}\delta\Sigma)$,
the natural gradient of $F$ is $\dot\mu=-\mathrm T^{-1}\Sigma\nabla^2U\mu$ and
$\dot\Sigma=-2\Sigma\bigl(\tfrac12\mathrm T^{-1}\nabla^2U-\tfrac12\Sigma^{-1}\bigr)\Sigma=\Sigma-\mathrm T^{-1}\Sigma\nabla^2U\Sigma$.
At $\Sigma=\mathrm T(\nabla^2U)^{-1}$ the second right side vanishes and the first is $-\mu$; the field
is smooth, so by uniqueness the slice is invariant and $\mu(t)=e^{-t}\mu(0)$. The residuals are linear
functionals of the mean that equal $-1$ at $\mathbf x^\circ$, which gives the stated residuals, and
$\lvert r^b_i\rvert=\lvert r^m_i\rvert$ excludes
$\lvert r^b_i\rvert\leq2\varepsilon_{\mathrm m}<\tfrac12\leq\lvert r^m_i\rvert$; the interval of
\cref{prop:obs-slow-model-link-dichotomy}(ii) is stated for $\varepsilon_{\mathrm m}\leq1/25$, where
$2\varepsilon_{\mathrm m}<\tfrac12$. \hfill$\square$

\noindent\emph{Proof of \Cref{prop:link-history-prior}.}
(a) Since $\mathscr S\geq0$, $\mathfrak U\geq0$, and $\mathbf m>0$, the exponent in
\eqref{eq:link-history-prior} is finite and at most $-\sum_n\mathscr S(\mathring\omega^{(n)})\leq0$.
The integrand is therefore positive everywhere and at most
$\prod_ne^{-\mathscr S(\mathring\omega^{(n)})}$, whose integral is $Z_{\mathscr S}^{n_1+1}\leq1$, so
$0<Z_{\mathrm{hist}}\leq Z_{\mathscr S}^{n_1+1}$ and the density bound follows. A Gaussian law
integrates $e^{\mathrm c_2\lVert\cdot\rVert}$.

(b) By Tonelli's theorem each $\mathring Z^{\leftarrow}_n$ is measurable; by induction from $n_1$, its integrand is
positive and at most $1$ and $\nu^{\mathscr T}$ is a probability law, so $\mathring Z^{\leftarrow}_n\in(0,1]$. Each
kernel of \eqref{eq:link-history-kernels} has total mass one by \eqref{eq:link-backward-functions},
and $\int e^{-\mathscr S}\mathring Z^{\leftarrow}_0d\nu^{\mathscr T}=Z_{\mathrm{hist}}$ by iterated integration of
\eqref{eq:link-history-prior}. At a history, the root density is
$Z_{\mathrm{hist}}^{-1}e^{-\mathscr S(\mathring\omega^{(0)})}\mathring Z^{\leftarrow}_0(\mathring\omega^{(0)})$ and
the kernel density of the step from $n$ to $n+1$ is
$e^{-\mathscr S(\mathring\omega^{(n+1)})-\mathbf m\mathfrak U(\mathring\omega^{(n+1)},\mathring\omega^{(n)})}$ times
$\mathring Z^{\leftarrow}_{n+1}(\mathring\omega^{(n+1)})/\mathring Z^{\leftarrow}_n(\mathring\omega^{(n)})$; the
ratios telescope, and since $\mathring Z^{\leftarrow}_{n_1}\equiv1$ the product is the density of
\eqref{eq:link-history-prior}.
So $\mathbb P^{\mathrm{hist}}$ is the law of this chain, the conditional law of
$\mathring\omega^{(n+1)}$ given $\mathring\omega^{(0:n)}$ is
$\mathbb P^{(n+1\leftarrow n)}_{\mathrm{hist}}(\cdot\given\mathring\omega^{(n)})$, and the kernels are the conditionals.
The chain that normalizes each increment factor separately has as kernels the factors
$e^{-\mathscr S(\mathring\omega')-\mathbf m\mathfrak U(\mathring\omega',\mathring\omega)}$ against
$\nu^{\mathscr T}(d\mathring\omega')$ divided by $\mathring Z^{\leftarrow}_{n_1-1}(\mathring\omega)$, which is
the normalizer of one increment factor, so the density of that chain relative to
$\mathbb P^{\mathrm{hist}}$ is proportional to
$\prod_{n<n_1}\mathring Z^{\leftarrow}_{n_1-1}(\mathring\omega^{(n)})^{-1}$, and by Fubini's theorem this product is
constant almost everywhere exactly when $\mathring Z^{\leftarrow}_{n_1-1}$ is. In the scalar case,
$\mathring Z^{\leftarrow}_0(\mathring\omega)=(1+\mathbf m)^{-1/2}\exp[-\mathbf m\mathring\omega^2/(2(1+\mathbf m))]$
is not constant; equivalently, the Gibbs law has precision matrix
$\left(\begin{smallmatrix}1+\mathbf m&-\mathbf m\\-\mathbf m&1+\mathbf m\end{smallmatrix}\right)$, so
its first marginal has precision $(1+2\mathbf m)/(1+\mathbf m)\neq1$, whereas the other chain starts
from $\nu^{\mathscr T}$, of precision $1$. For the
latent route, the joint
$\mathbb P^{\mathrm{hist}}(d\mathring\omega^{(0:n_1)})\prod_n\mathbb P^{\mathrm{ag}}_n(dy^{(n)}\given\mathring\omega^{(n)},X)\prod_n\mathbb P^{\mathrm{obs}}(do_n\given\mathbf Y^{(n)},X)$,
with the chain factorization substituted, is \eqref{eq:full-graph-episode-joint} for the stated root,
transitions, and readout: each transition is a Markov kernel from
$Z^{(n-1)}_{\mathrm{par}}=\mathring\omega^{(n-1)}$ to $\mathbf Y^{(n)}$ composed of two Markov kernels,
and the root is a law of $\mathbf Y^{(0)}$. Every kernel reads structural data and earlier latent
coordinates, as \cref{req:gen-typing-prohibition} requires, and \cref{prop:full-graph-episode-joint}
gives the listed consequences, the predictive prior being
\eqref{eq:full-graph-episode-predictive} with $Z^{(n)}_{\mathrm{par}}=\mathring\omega^{(n)}$.

(c) With a static frame $U_i$, the twist of the identity link of agent $i$ is
$U_i^{-1}e_GU_i=e_G$ by \eqref{eq:geo-edge-twist}, independent of every variable, so it is the
identity twist of each time-like step in \eqref{eq:link-timelike-holonomy}, and
\eqref{eq:link-history-prior} contains no other link. A passive re-expression leaves every twist
unchanged (\Cref{prop:geo-edge-twist-invariance}), so it neither produces nor removes this identity
twist. \hfill$\square$

\noindent\emph{Proof of \Cref{prop:link-map-inertia}.}
(a) Differentiating \eqref{eq:link-history-action} in $\mathring\omega^{(n)}$ gives the equations,
since $\mathring\omega^{(n)}$ enters the increments to $n-1$ and $n+1$ in the interior and one increment
at each end. Under the growth condition $\mathscr S^{\mathrm{hist}}$ is continuous and coercive on a finite-dimensional
space, so it attains its minimum.

(b) By \cref{lem:link-exp-derivative}(a),
$\tfrac d{dt}e^{\mathbf B+t\mathbf T}\vert_{t=0}=\mathbf L(\mathbf B,\mathbf T)$. For each $t\in[0,1]$
the map $\mathbf X\mapsto e^{t\mathbf B}\mathbf Xe^{(1-t)\mathbf B}$ has the adjoint
$\mathbf X\mapsto e^{t\mathbf B^*}\mathbf Xe^{(1-t)\mathbf B^*}$ for $\langle\cdot,\cdot\rangle_{\mathsf G}$,
because $(e^{t\mathbf B})^*=e^{t\mathbf B^*}$, so the adjoint of $\mathbf L(\mathbf B,\cdot)$ is
$\mathbf L(\mathbf B^*,\cdot)$. The differential of $\mathfrak Q$ at $\mathbf H$ is
$\mathbf T\mapsto\langle\mathbf H-I,\mathbf T\rangle_{\mathsf G}$, and the differential of
$\rho(\mathring\Theta^{\leftarrow}_e)$ in the direction $\delta\mathring\omega_e$ is
$e^{-\mathrm d\rho\mathring\omega'_e}\mathbf L(\mathrm d\rho\mathring\omega_e,\mathrm d\rho\delta\mathring\omega_e)$ in the
earlier argument and
$-\mathbf L(-\mathrm d\rho\mathring\omega'_e,\mathrm d\rho\delta\mathring\omega_e)e^{\mathrm d\rho\mathring\omega_e}$ in the
later one; moving the factors across the inner product and projecting onto $\mathrm d\rho(\mathfrak g)$
gives \eqref{eq:link-timelike-gradients}, and \eqref{eq:link-map-inertia-defect} is stationarity of
\eqref{eq:link-history-action}. The function $\mathfrak U$ is smooth, nonnegative, and zero on the
diagonal, so its gradient vanishes there; and
$\rho(\mathring\Theta^{\leftarrow}_e)=I+\mathrm d\rho(\mathring\omega_e-\mathring\omega'_e)+O(\lVert(\mathring\omega,\mathring\omega')\rVert^2)$
gives $\mathfrak U=\tfrac12\lVert\mathring\omega-\mathring\omega'\rVert^2+O(\lVert(\mathring\omega,\mathring\omega')\rVert^3)$,
so the mixed second derivative at the origin is $-I$. Let
$\mathrm R_1(\mathring\omega',\mathring\omega)=\partial_{\mathring\omega}\mathfrak U-(\mathring\omega-\mathring\omega')$.
Then $\mathrm R_1(\mathring\omega,\mathring\omega)=0$ and
$\mathrm R_1(\mathring\omega',\mathring\omega)=\int_0^1D_{\mathring\omega'}\mathrm R_1(\mathring\omega+t(\mathring\omega'-\mathring\omega),\mathring\omega)[\mathring\omega'-\mathring\omega]dt$,
where $D_{\mathring\omega'}\mathrm R_1$ vanishes at the origin and has derivative bounded by
$\sup\lVert D^3\mathfrak U\rVert$ on the convex set
$\{\lVert\mathring\omega\rVert,\lVert\mathring\omega'\rVert\leq\mathrm r\}$; since
$\lVert(\mathring\omega,\mathring\omega')\rVert\leq\sqrt2\mathrm r$ there, the bound follows, and the same
argument applies to $\partial_{\mathring\omega'}\mathfrak U-(\mathring\omega'-\mathring\omega)$. For the
quadratic form at a given $\mathring\omega$, with $\mathring\omega'=\mathring\omega+\delta\mathring\omega$,
the differential of the later argument gives
$\rho(\mathring\Theta^{\leftarrow}_e)=I-\mathbf L(-\mathrm d\rho\mathring\omega_e,\mathrm d\rho\delta\mathring\omega_e)e^{\mathrm d\rho\mathring\omega_e}+O(\lVert\delta\mathring\omega\rVert^2)$,
and $\mathbf L(0,\mathbf T)e^{0}=\mathbf T$ at the origin.

(c) Write $\mathrm E(\mathring\omega)=\tfrac{\mathbf m}2\sum_{n<n_1}(\mathring\omega^{(n+1)}-\mathring\omega^{(n)})^2+\sum_{0<n<n_1}\mathfrak V(\mathring\omega^{(n)})$
for histories with the given ends. Replacing every value by its clamp to $[\mathrm a,\mathrm b]$ does
not increase $\mathrm E$: clamping is $1$-Lipschitz, and by the hypothesis on $\mathfrak V$ outside
$[\mathrm a,\mathrm b]$ it does not raise the potential terms. If some value leaves $[\mathrm a,\mathrm b]$, the
step that enters a maximal run of such values from a value in $[\mathrm a,\mathrm b]$ strictly
shortens, so the decrease is strict. Hence minimizers over $[\mathrm a,\mathrm b]^{n_1-1}$, which
exist by compactness and continuity, are global, and every global minimizer takes its values in
$[\mathrm a,\mathrm b]$. Now let $\mathrm z_0\leq\cdots\leq\mathrm z_{n_1}$ be the values of such a
history in increasing order, so $\mathrm z_0=\mathrm a$ and $\mathrm z_{n_1}=\mathrm b$, and let
$\mathrm g_{\mathrm l}=\mathrm z_{\mathrm l+1}-\mathrm z_{\mathrm l}$. No value lies strictly inside a gap, so each step from
$\mathring\omega^{(n)}$ to $\mathring\omega^{(n+1)}$ covers a set of whole gaps and its square is at
least the sum of their squares. A history from $\mathrm a$ to $\mathrm b$ crosses each gap of
positive length an odd number of times, so summing over steps gives
$\sum_n(\mathring\omega^{(n+1)}-\mathring\omega^{(n)})^2\geq\sum_{\mathrm l}\mathrm g_{\mathrm l}^2$, the kinetic sum of
the increasing rearrangement; a decreasing step crosses some gap of positive length downward, which
forces at least three crossings of it and a strict inequality. The rearrangement keeps $\mathrm a$
and $\mathrm b$ at the ends and the multiset of interior values, so it leaves the potential sum
unchanged. A history that is not nondecreasing therefore has a strictly smaller competitor, and
every most probable history is nondecreasing. For strictness, let $[\mathrm p,\mathrm q]$ be a
maximal run of equal values $\mathrm x$ with $\mathrm q\geq\mathrm p+1$. If $0<\mathrm p$ and
$\mathrm q<n_1$, \eqref{eq:link-map-inertia} at $\mathrm p$ and at $\mathrm q$ gives
$\mathbf m(\mathrm x-\mathring\omega^{(\mathrm p-1)})=-\mathfrak V'(\mathrm x)=\mathbf m(\mathrm x-\mathring\omega^{(\mathrm q+1)})$,
which contradicts $\mathring\omega^{(\mathrm p-1)}<\mathrm x<\mathring\omega^{(\mathrm q+1)}$. If
$\mathrm p=0$, then $\mathrm x=\mathrm a$, $\mathrm q<n_1$, and \eqref{eq:link-map-inertia} at
$\mathrm q$ with $\mathfrak V'(\mathrm a)=0$ gives $\mathring\omega^{(\mathrm q+1)}=\mathrm a$, against
maximality; the case $\mathrm q=n_1$ is symmetric, and $\mathrm p=0$ with $\mathrm q=n_1$ would give
$\mathrm a=\mathrm b$. \hfill$\square$

\noindent\emph{Proof of \Cref{prop:link-linearized-equation}.}
(a) The anchor term has gradient $0$ and Hessian $\mathring\Sigma^{-1}$ at $0$. Each $\mathfrak Q_p$ is
smooth, a polynomial in the entries of a product of matrix exponentials, and
$\mathfrak Q_p(\mathring\omega)=\tfrac12\lVert\mathbf C_p\mathring\omega\rVert_{\mathsf G}^2+O(\lVert\mathring\omega\rVert^3)$
by \cref{prop:link-loop-defect}(e), so by uniqueness of Taylor coefficients its gradient vanishes at
$0$ and its Hessian there is $\mathbf C_p^\top\mathbf C_p$. The Gaussian increment is the quadratic form
of $\mathbf m(-\Delta_n)$, and the time-like defect has the same second-order part by the expansion in
the proof of \cref{prop:link-map-inertia}(b), with vanishing gradient at $0$. For a record $a$ with
scope $\{i,j\}$ on $e:j\to i$ and $\mathring\Omega=\mathrm d\rho\mathring\omega_e$, independence of $k_i$
and $k_j$ gives
$\E[E_{a,o}]=\tfrac12[\lVert\mu_i-e^{\mathring\Omega}\mu_j\rVert_{W_a}^2+\Tr(W_a\Sigma_i)+\Tr(W_ae^{\mathring\Omega}\Sigma_j(e^{\mathring\Omega})^\top)]$,
dropping the episode index. With $e^{\mathring\Omega}=I+\mathring\Omega+\tfrac12\mathring\Omega^2+O(\lVert\mathring\Omega\rVert^3)$,
the first-order part is
$\Tr(W_a\mathring\Omega\Sigma_j)-(\mu_i-\mu_j)^\top W_a\mathring\Omega\mu_j=-\Tr(\mathbf J_a^\top\mathring\Omega)$ and the
second-order part is
$\tfrac12[\Tr(W_a\mathring\Omega\E[k_jk_j^\top]\mathring\Omega^\top)+\Tr(W_a\mathring\Omega^2\Sigma_j)-(\mu_i-\mu_j)^\top W_a\mathring\Omega^2\mu_j]=\tfrac12[\Tr(W_a\mathring\Omega\E[k_jk_j^\top]\mathring\Omega^\top)-\Tr(\mathbf J_a^\top\mathring\Omega^2)]$,
because $\mathbf J_a^\top=(\mu_j(\mu_i-\mu_j)^\top-\Sigma_j)W_a$. For a record on a reversed copy,
$\mathring\Omega=-\mathrm d\rho\mathring\omega_{\bar e}$ with $\bar e\in E_+$, which changes the sign of the
first-order part only. Summing with the weights gives $\nabla\mathscr S^{\mathrm{hist}}(0)=-\mathbf j$ and
$\nabla^2\mathscr S^{\mathrm{hist}}(0)=\mathbf K+\mathbf N_0$.

(b) The loop, episode, and matter forms are nonnegative, the last because
$\Tr(W_a\mathring\Omega\E[k_jk_j^\top]\mathring\Omega^\top)=\E\lVert W_a^{1/2}\mathring\Omega k_j\rVert^2$, so
$\langle\mathring\omega,\mathbf K\mathring\omega\rangle\geq\mathrm c_{\mathring\Sigma}\lVert\mathring\omega\rVert^2$,
and the bounds follow from $\lVert\mathbf K^{-1}\rVert\leq1/\mathrm c_{\mathring\Sigma}$.

(c) A represented generator $\mathring\Omega$ has
$\lVert\mathring\Omega\rVert_{\mathsf G}=\lVert(\mathsf G^{\mathrm{fib}})^{1/2}\mathring\Omega(\mathsf G^{\mathrm{fib}})^{-1/2}\rVert_{\mathrm F}$,
so that
$\lVert\mathring\Omega\rVert_{\mathrm F}^2\leq\operatorname{cond}(\mathsf G^{\mathrm{fib}})\lVert\mathring\Omega\rVert_{\mathsf G}^2$.
Then $\lvert\Tr(\mathbf J_a^\top\mathring\Omega^2)\rvert\leq\lVert\mathbf J_a\rVert_{\mathrm F}\lVert\mathring\Omega\rVert_{\mathrm F}^2$
and $\lvert\Tr(\mathbf J_a^\top\mathring\Omega)\rvert\leq\lVert\mathbf J_a\rVert_{\mathrm F}\lVert\mathring\Omega\rVert_{\mathrm F}$,
summed over records and combined with the Cauchy--Schwarz inequality, give the bounds on
$\mathrm c_0$ and $\lVert\mathbf j\rVert$, the first because $\mathbf N_0$ is self-adjoint. Write
$\nabla\mathscr S^{\mathrm{hist}}(\mathring\omega)=-\mathbf j+\mathbf K\mathring\omega+\nabla\mathrm R(\mathring\omega)$, so that
$\nabla\mathrm R(0)=0$, $\nabla^2\mathrm R(0)=\mathbf N_0$, and
$\lVert\nabla^2\mathrm R(\mathring\omega)\rVert\leq\mathrm c_0+\mathrm c_{\mathrm r}\lVert\mathring\omega\rVert$ for
$\lVert\mathring\omega\rVert\leq\mathrm r$. The critical points of $\mathscr S^{\mathrm{hist}}$ are the fixed points of
the map $\mathring\omega\mapsto\mathbf K^{-1}(\mathbf j-\nabla\mathrm R(\mathring\omega))$. On the ball of radius
$\mathrm r_1=2\lVert\mathbf j\rVert/\mathrm c_{\mathring\Sigma}\leq\mathrm r/2$,
$\lVert\nabla\mathrm R(\mathring\omega)\rVert\leq(\mathrm c_0+\mathrm c_{\mathrm r}\mathrm r)\lVert\mathring\omega\rVert\leq\tfrac12\mathrm c_{\mathring\Sigma}\mathrm r_1$,
so the map sends this ball into itself, its values having norm at most
$\lVert\mathbf j\rVert/\mathrm c_{\mathring\Sigma}+\mathrm r_1/2=\mathrm r_1$, and it has Lipschitz constant at
most $(\mathrm c_0+\mathrm c_{\mathrm r}\mathrm r)/\mathrm c_{\mathring\Sigma}\leq\tfrac12$ there, so it has exactly
one fixed point $\mathring\omega_\star$ in that ball. On the ball of radius $\mathrm r$,
$\nabla^2\mathscr S^{\mathrm{hist}}\succeq(\mathrm c_{\mathring\Sigma}-\mathrm c_0-\mathrm c_{\mathrm r}\mathrm r)I\succeq\tfrac12\mathrm c_{\mathring\Sigma}I$,
so $\mathscr S^{\mathrm{hist}}$ is strongly convex on that convex set and has there at most one critical point, which is
its unique minimizer on the set. Finally
$\mathring\omega_\star-\mathring\omega_{\mathrm{lin}}=-\mathbf K^{-1}\nabla\mathrm R(\mathring\omega_\star)$ and
$\lVert\nabla\mathrm R(\mathring\omega_\star)\rVert\leq(\mathrm c_0+\tfrac12\mathrm c_{\mathrm r}\lVert\mathring\omega_\star\rVert)\lVert\mathring\omega_\star\rVert$,
which with $\lVert\mathring\omega_\star\rVert\leq\mathrm r_1$ gives \eqref{eq:link-linearization-error}. Since
$\mathrm c_0$ and $\lVert\mathbf j\rVert$ are bounded by multiples of the currents, the right side is
of second order in them. \hfill$\square$
